\section{User's guide}

\subsection{General concept}

As discussed in Section~3 the set of transport equations for a tokamak
varies depending on a problem to be solved.
The system \req{main} can be either shortened or appended with additional
equations.
Right hand sides of the system can also have quite different forms for
different applications.
However, the uncertainty in the transport matrix \req{fluxes} is
incomparably higher.
It would not be an overstatement to say that there are hundreds of
choices for transport coefficients in the matrix \req{fluxes}.
Therefore, to be efficient the transport simulation of a tokamak
requires a convenient and flexible tool for creating numerical codes
based on different transport models.
The programming system Astra was designed for this purpose.

Proceeding from this requirement the Astra code considers the set
\eqs\ref{main})-(\ref{fluxes}) as a pattern which can be filled in or
left empty according to user's request.
Astra is not a numerical program in the conventional sense.
A pre-determined program exists for the service parts of the Astra
system only. 
The transport kernel of the code is created every time by a code builder
and includes only those equations,
transport coefficients, sources, sinks and additional modules that were
explicitly requested by the user in the description of transport task.
It does not mean, however, that any change requires re-building the
program. 
On the contrary, the code is designed for an interactive work and
provides a lot of techniques for computation control and correction.
In particular, the interactive mode of operation is convenient for an
interpretative modeling or for analysis of experimental data.

Developing the concept of the code we tried to meet two contradictory
requirements. 
The first one is minimizing user's efforts for learning the code and for
starting any reasonable simulation. 
This implies that the user has to take only those decisions and supply
only this information which cannot be obtained elsewhere.
For instance, the type of device (tokamak or stellarator), its
geometrical size, magnetic field, plasma current have to be specified by
the user.
It is also exclusively left at user's discretion to select a way of data
treatment and data presentation.
The second requirement is to avoid any hidden assumptions which can
result in unexpected or difficult for understanding results and thus
make usage of the code even more complicated if possible at all.  
As a compromise, we adopted a number of default transformations which
are in effect only if no explicit request is provided by the user.

The following questions should be answered before creating any transport
code:
\begin{description}
\item[\ \ 1.] What subset of \eqs{\ref{GSp}}), (\ref{main}), (\ref{aux}) is
to be solved?
\item[\ \ 2.] What boundary conditions should be imposed?
\item[\ \ 3.] What quantities and in which form should appear as an output?
\item[\ \ 4.] What parameters (major and minor radii of a torus, magnetic
field, plasma current, etc.) and what initial conditions should be used?
\end{description}
\nopagebreak
In the Astra system the fourth question is separated from the others.
It is reasonable because the first three items specify the simulation
model, i.e. a set of equations to be solved and a form of presentation
of the results, whereas the forth specifies where to apply the
simulation model. 
By providing the answers to the first three questions the user completely
defines how he wants to perform a simulation.
Then Astra automatically builds the source code for the specified transport
problem, compiles and loads it.
We will call the generated program code a transport model or simply a
model.

The answer to the fourth question fills a model with machine and shot
specific contents.
Although the machine and plasma parameters can also be included in a
model it is reasonable not to do this.
Then the same simulation model can be applied to different devices and
different parameters that are read from a data file when the program
starts execution. 
In addition to the global device parameters and the initial data for
transport simulation, this file can also contain the full available
information of plasma parameters, e.g. the time evolution of plasma
current, plasma density $n_e$, the electron $T_e$ and ion $T_i$
temperatures, plasma position, $Z_{eff}$, impurity densities, radiated
power, auxiliary heating power, driven current and many others. 
Essentially, it stores the whole set of experimental data for a specific
shot.
In addition, database shot record may include results of calculation from
other codes thus providing implicit link between Astra and other codes,
when external calculations may be made in advance.
All this informations along with specifying discharge parameters can be used
in the interpretative transport modeling for comparison with the simulation
results in order to refine and develop the transport model.
In the analysis mode, the experimental information is used for different
kinds of data analysis and processing.

Thus, {\bf the model} is a transport code created for the system of
transport equations \req{main} with specified transport coefficients
\req{fluxes}, initial and boundary conditions.
It can include also a variety of other equations (e.g. different current
drive and heating schemes, cold neutrals, SOL model, impurity diffusion),
formulae and all possible combinations and processing the obtained data. 
In addition to the transport analysis, the Astra system provides a
powerful apparatus for data processing and presentation.

{\bf The data file} is a set of time and radially dependent data
related to a specific shot in a given device and containing available
information about the shot.
The data from the record are read by a model and can be used as the
initial and boundary conditions, for comparison of stored experimental
data with the results of simulation, for different treatments of
experimental data and all other types of data analysis.

A number of different models and data records exist in the Astra system
independently.
The rules for the creation of a model are described in
Section ~5.2
The data file format is described in Section~5.4.

%\vfill\eject

\subsection{Setting a model}

\subsubsection{Model file}

The main step in creation of a simulation code within the Astra
system is a formulation of the transport problem to be solved and
selection of the form for presentation of the results.
This formulation is put into a special file where all requests
for transport simulation, data manipulation and result presentation 
are specified using Astra model description language.
After the code starts the file is submitted to the Astra compiler 
which processes all the instructions contained in the file and creates
first a Fortran/C program and then the executable transport code.
In most cases, these steps remain hidden for the Astra user who deals
exclusively with the model or prototype file. 
In what follows, we will not make a distinction between the transport
code as such and the prototype file.
Moreover, for brevity the prototype file will be also called
{\bf a model file} or just a model. 

Each model file contains a comprehensive and compact (typical length is
3~kB) information about a corresponding transport code (approximate
length of all source files 3 MB, binary executable 10~MB) created by the
Astra system. 
Because the model file is a unique image of the transport code, only
these files need to be retained in order to repeat any simulation when
such a necessity arises.
All later Astra versions maintain compatibility with model files of
preceding versions. 
The user can have simultaneously a number of model files which are all
stored in the directory \tw{AWD/equ/}.\footnote
{See footnote in the page \pageref{AWD}.} 
Few syntax rules are adopted in order to present user's request for
a transport code in the most clear, visible and concise form.

\subsubsection{General rules}

A model file is a text (ASCII) file.
Although a length of this file is restricted the restriction practically
never occurs. 
All instructions in a model are case insensitive.
Maximal length of each line in the file is 132 symbols including spaces
and tabulations which, however, are ignored and can be used to make a
file more readable.
The last line in the file should have a symbol \tw{<CR>} (carriage
return) on the end otherwise it is ignored  during processing.

Now we define a set of symbols, which have a special meaning
and serve to describe Astra instructions or commands. 
Those are
\begin{equation}\label{specsym}
\mbox{\tt :\qq =\qq\bs \qq\symbol{95}\qq ;\qq !}
\end{equation}
Every line in a model is considered as a command if anyone of the four
first symbols in \req{specsym}, i.e. 
{\tt ':', '=', '\bs ', '\symbol{95}'~} 
is present in this line.
Several commands can be joined in one line if they are separated 
from one another with a semicolon {\tt ';'}.

All other lines are ignored by the Astra compiler and, therefore, they
can be used as commentaries.
A command line can also be converted to a comment if the first symbol
in it is an exclamation sign \tw{'!'}.
In other words, the sign \tw{'!'} switches off all special symbols on 
the right of itself until the end of line (\tw{<CR>}) or the nearest
\tw{';'}. 
Additionally, if the 1-st symbol in a line is \tw{'!'} then the whole
line is ignored.

\vskip 1ex\noindent
{\bf Example:}\\
{\tt\hspace*{1mm}
~~~~~~~~~~~~~Boundary condition for the density\smallskip\\
!NEB=NEX;~~QNB=100;~~QNNB=30;
~This line is switched off by the leading '!'\smallskip\\
\symbol{32}!~NEB=3;~~QNB=500;~~!QNNB=700;
~Here the second instruction is effective 
}\bigskip\\
This example contains only one command: \tw{QNB=500}.
All other instructions are suppressed.

\paragraph{Order of operators.}

When creating an executable code the Astra compiler analyzes each
command in a model file and then creates a sequence of corresponding
Fortran operators.
(Sometimes, several operators correspond to one command in a model file.)
During a compilation all commands are split in classes.
The first class consists of assignment statements (Section 5.2.3) which
have a scalar variable on the left hand side.
The second class includes assignments of vector variables.
Finally, the third class is composed of external modules (plug-in
modules, Section 5.2.4) and transport equations (Section 5.2.5).  
These three classes are treated after one another.
With a few exceptions, an order of operators in each class is retained
the same as in the original model file. 

\subsubsection{Arithmetic expressions}

Arithmetic expressions in the Astra code basically obey the Fortran
syntax. 
Operands of an expression can be constants, scalar and vector variables,
Astra formulae and functions. 
All arithmetic operations and Fortran intrinsic functions can be
used in the expressions. 

As in Fortran, arithmetic expressions appear mainly in {\bf assignment
statements} i.e. those command lines which include the sign \twd{'='}.
On the left of assignment statements can be the Astra scalars and
vectors only. 
If the left hand side of the assignment statement is a scalar than the
right hand side of it must be a scalar otherwise an error is reported.
If the left hand side is a vector than the right hand side can be either
a scalar or a vector.

\vskip 1ex\noindent
{\bf Examples:}\smallskip\\ %{\rm(1)}\\
{\tt\hspace*{1mm}
~~~~~~Power source of 1 MW with the Gaussian distribution in radius
\smallskip\\
CAR1=exp(((RHO/ROC-0.5)/0.1)**2);~~~~~~~~CF1=VINT(CAR1B);
\smallskip\\ \hspace*{1mm}
~~~~~~~~~~~~~~~~~~~~~and a rectangular modulation in time (100Hz)
\smallskip\\
QECR=anint(.1*sin(100*GP2*TIME)+.5);~~PEX=QECR*CAR1/CF1;
}\bigskip\\
This example illustrates using expressions in a model. 
Argument of the exponent in the first command shows that constants,
scalars and vectors can be used on the right of assignment statement.
The result is a vector.
The second command calculates the scalar
{\tt CF1} so that {\tt CAR1/CF1} is normalized to 1 MW.
The third command uses two intrinsic Fortran functions and, finally,
the forth one assigns the vector variable {\tt PEX} that later can be
used as heating source in a transport equation.

The next example shows that predefined Astra expressions can be used in
a model to construct quite complicated quantities.
{\tt%{\rm(2)}
\smallskip\\
\hspace*{1mm}
~~~~~~~Simplified critical gradient model
\smallskip\\
CAR1=RTOR*grad(TI)/TI-RLTCR;~~~~~~CAR1=sqrt(step(CAR1)*CAR1);
\smallskip\\
CAR1=CAR1-CV1*ROTSH/GITG0;~~~~~~~~XI=HNCHI+CHE1*step(CAR1)*CAR1;
}\medskip\\
Here the following model for the ion heat conductivity is implemented
$$
\begin{array}{c}\ds
\eta_i=\frac{R_0 T_i'}{T_i}
,\qq
\gamma^{ITG}={\rm max}\lb(0,\gamma_0\sqrt{\eta_i-\eta_{i,cr}}\rb)
,\qq
\chi_i^{ITG}
=\frac{\chi_0}{\gamma_0}\,
{\rm max}\,\lb(0,\gamma^{ITG}-\omega_{E\x B}\rb)
,\bigskip\\ \ds
\chi_i=\chi_i^{PS}+\chi_i^{neo}+\chi_i^{ITG}
\end{array}
$$
and the following quantities from the Astra library (Astra formulae) are
involved\\  
\hspace*{3mm}
{\tt RLTCR~} $\eta_{i,cr}=\lb(\frac{R_0\na T_i}{T_i}\rb)_{cr},$ 
critical ion temperature gradient (ITG),\\
\hspace*{3mm}
{\tt ROTSH~} $\omega_{E\x
B}=\frac{B_p}{B_0}\lb(\frac{E_\rho}{B_p}\rb)',$\label{ROTSH} 
rotational shear,\\
\hspace*{3mm}
{\tt GITG0~} $=\gamma_0$ is related to the
linear increment of the ITG instability $\gamma^{ITG}$,\\
\hspace*{3mm}
{\tt HNCHI~} $\chi_i^{PS}+\chi_i^{neo}$ is the
neoclassical heat conductivity \cite{CH}.\\ 
Two free parameters are included in the model:
the stiffness $\chi_0=\tt CHE1$ 
that characterizes transport properties of the nonlinear phase of the
instability and the parameter {\tt CV1} which enables switching on and
off transport (or
instability) suppression due to the plasma rotation. 

To summarize the interpretation rules for Astra expressions we say that
a type of the result is determined by the left hand side of the
assignment statement. 
An assignment of a scalar to a vector variable is not considered as an
error so that the vector on the left retaining its type is viewed as
radially independent.
If a vector is attempted to be assigned to a scalar then a warning
message is reported although the error is not considered as severe.

\paragraph{Vector variables at given radii.}

As discussed, the Astra compiler discriminates different types
of variables and allows using them in common context without explicit 
indication of the variable type.
However, sometimes a need arises to evaluate vector quantities at
given radial positions. 
This can be done in the conventional way.
For instance, \tw{TE(0)} will be understood as the electron temperature
at the magnetic axis $T_e(0),$ \tw{TE(ABC)} as the electron temperature
at the plasma edge and \tw{~TE(0.3)} as the electron temperature at the
radius $a=$0.3 m.
More complicated constructions are also allowed.
For instance, \tw{TE(AFVAL(MU,0.5))} gives a value of $T_e$ at the
resonance surface $q=2.$

A reference to the coordinate $a$ is selected as most convenient for 
associating it with the laboratory reference frame because this
coordinate can be thought of as the minor radius in the mid-plane
(see \req{3mom}).
However, Astra vectors can be calculated also as functions of other
flux coordinates using auxiliary functions described in Section~4.8.3. 
So \tw{~NI(AFR(0.3))} gives the ion density $n_i(\rho=0.3\rm m),$
and \tw{~NI(AFX(0.3))} at $\rho_N=\rho/\rho_B=0.3.$
One can also use \tw{~NI(0.3*ABC)} to obtain $n_i$ at $a_N=a/a_B=0.3$ 
or, what is the same, at $a=0.3\,a_B=0.3\tt*ABC.$ 

Shortcuts are provided for vector quantities evaluated at the magnetic
axis and at the plasma edge. 
Namely, adding {\tt~C~} to a vector variable name one obtains its
value at the centre and adding {\tt~B~} at the edge
\footnote{One exception of this rule is discussed at the end of 
Section~5.2.5}.
It means that \twd{TEC} stands for \twd{TE(0),} \twd{MUC} for 
\twd{MU(0),} \twd{TEB} for 
${\tt TE}|_{a=a_B}={\tt TE}|_{\rho=\rho_B}=\tt TE(ROC),$ 
\twd{MUB} for \twd{MU(ROC)} and so on.

Astra expressions which are implemented as functions (but not as
formulae) can also be calculated at given radii according to the same
syntax rules.
Thus \twd{NEAVB} yields the electron density averaged over the whole 
plasma volume while \twd{NECHC} gives the line averaged density along 
the central chord.
Finally, one can use two Astra functions \twd{ATR} and \twd{ATX} (see 
Section~4.8.3) in order to evaluate Astra vectors as functions of
the variables $\rho$ and $\rho_N,$ respectively.

\subsubsection{Connection of plug-in modules}

The next subject to be discussed is a connection of auxiliary
packages to the Astra code. 
On one hand, development of external packages provides the most
powerful way to include various physical processes in the code. 
In this way the user can also get an access to almost any part of the
code and make nearly any 
changes and additions to the resulting simulation model.
On the other hand, it is completely at user's discretion to trace
possible consequences of such an interference in the code flow.
Many plug-in modules are already available and were discussed in the
previous section.
Others can be developed using one or few of the existing modules as
templates. 

Front end of any plug-in module is a subroutine called by Astra.
Their inclusion in a simulation model is implemented through the
instruction of the type\footnote
{Here and below the braces \{...\} in a command description
enclose an optional part of the command.}.
\begin{equation}\label{sbr}
\mbox{\tt Name(\it Argument list}\mbox{\tt):\symbol{123}}\Delta_t
\tt:t_{start}\tt:t_{end}\mbox{\tt:Letter\symbol{125}}
\end{equation}
Only the first colon `:' is required in this line all others can be
omitted together with the corresponding control parameters.
However, if the $i\!$-th $(i=1,2,3,4)$ parameter is present it
should be preceded with $i$ colons.
This command results in calling a subroutine with the name {\tt "Name"}.
An argument list of the subroutine (if present) is directly converted
into a Fortran line without any transformations which are normally
performed with the Astra expressions. 
Therefore, the actual arguments can be constants, variable names, 
or expressions including constants and scalar variables only.
For instance, the instruction 
\\
{\tt\hspace*{3em}NEUT:;} 
\\
corresponds to the Fortran
statement\\ 
{\tt
%\symbol{32}\symbol{32}\symbol{32}\symbol{32}\symbol{32}\symbol{32}
\hspace*{3em}call~NEUT}\\
which calls the subroutine {\tt NEUT} calculating the neutral density 
{\tt NN} and temperature {\tt TN} (see Section 4.9.4).

In the last example, the subroutine will be called at each time step of
the simulation run. 
This can be changed by involving control parameters in the command
line \req{sbr}.
They have the following meaning
\bigskip\\
\begin{tabular}{|c|l|c|}\hline
\bf Parameter	&\bf Significance &\bf Default value	\\
\hline
$\tt \Delta_t$	& Time interval between calls	& 0	\\
$\tt t_{start}$& First time of calling	&$-\infty$\\
$\tt t_{end}$& Last time of calling	&$\infty$\\
{\tt Letter} 	& Control key			& None	\\
\hline
\end{tabular}
\bigskip\\
If $\tt\Delta_t$ is not specified a subroutine is called at each time
step. 
Parameters $\tt t_{start}$ and $\tt t_{end}$ define start and end
time 
when the module is activated. No calls are made beyond the specified time
interval $\tt t_{start}\leq{\it t}\leq t_{end}$.
If one of these parameters is not specified then corresponding conditions
do not apply.
Parameter {\tt Letter} is used to provide manual activation of the module
in the interactive mode of operation.
So, the instruction 
\\
{\tt\hspace*{3em}NEUT:0.1:::N;} 
\\
results in calling \tw{NEUT} every 100 ms and additionally each time
when the key combination ${\tt<Ctrl>N}$ is pressed. 
The control key {\tt Letter} can acquire any capital alphabetic value
except for \tw{'M'}.

\subsubsection{Transport equations}

As described in Section 4, the Astra code can solve up to seven transport
equations for plasma density, electron and ion temperatures, current
density and three unspecified functions.
Each equation can be involved by an explicit request only.
Alternatively, a rule can be provided for calculating any of those
quantities as functions of radius and time.
A type of evolution of a variable is specified by the
instruction of the form
\begin{equation}\label{Eq_com}
\mbox{\tt Name:\symbol{123}Type\symbol{125}\symbol{123}[Tag,Expr1]\symbol{125}\symbol{123}:Expr2\symbol{125}}
\end{equation}
where every term after the first \tw{':'} can be omitted. 
\begin{description}
\vspace*{-1ex}
\item[{\tt Name}] is one of the variables on the list: 
${\tt NE,~TE,~TI,~CU,~F1,~F2,~F3}.$
\vspace*{-1ex}
\item[{\tt Type}] is one of the two qualifiers \twd{Equation} or
\twd{Assignment} (each can be truncated up to one letter), default value is
{\tt~Type=Equation}. 
\vspace*{-1ex}
\item[{\tt[Tag,Expr1]}] Not applicable if \tw{~Name=CU}.
Otherwise, \tw{Expr1} defines the right boundary of the
interval where the transport equation will be solved.
An integer \tw{Tag} specifies a type of the radial coordinate defined by
the expression \twd{Expr1} \\
\hspace*{2em}{\tt Tag=0~} dimensional coordinate $a$ (may be omitted),\\
\hspace*{2em}{\tt Tag=1~} dimensionless coordinate $a/a_B,$\\
\hspace*{2em}{\tt Tag=2~} dimensionless coordinate $\rho_N=\rho/\rho_B.$
\vspace*{-1ex}
\item[{\tt Expr2}] is another expression that can be used if {\tt Name}
is one of ~{\tt TE}~~or~~{\tt TI} only.
The expression then determines $\gamma_e$ or $\gamma_i$ in \req{Amain}.
By default, $\gamma_e=\gamma_i=0.$
\end{description}
\vspace*{-1ex}
If the qualifier \tw{Type} is omitted then it is assumed to take the
value \tw{Type=Equation}. 
In case \tw{Type=Assignment}, the rest of the command string is ignored.
The new feature introduced in the version 5.3 is a possibility to set the
different right boundaries for different transport equations.
However, some restrictions are imposed.
\begin{enumerate}
\vspace*{-1ex}
\item If the right boundary is not given then it is set to
$\rho_B=\tt ROC$ (or $a_B=\tt ABC$). 
This maintains compatibility with the previous versions. 
\vspace*{-1ex}
\item For the poloidal flux (i.e. for \tw{Name=CU}) the right
hand boundary cannot be shifted and is always $\rho_B.$ 
\vspace*{-1ex}
\item For all other equations it cannot exceed the value $\rho_B.$
\vspace*{-1ex}
\item Unlike the outermost boundary position, $\rho_B,$ which can be
located arbitrary (continuously varying) on the difference grid, all
``internal'' 
boundaries are attributed to the nearest grid point in $\rho$ (jumping).
\end{enumerate}
\vspace*{-1ex}
When a model is converted into a Fortran code, the transport equations
and the plug-in modules are put in the same order as they are requested
in the model. 
Exception is made for the subroutines \tw{MIXINT} and \tw{TSCTRL} which
are always placed after all transport equations.
Although, in most cases, the order of transport equations is not
essential, sometimes, one has to take this convention into account.

The following example shows setting transport equation for the plasma
density 
\begin{equation}\label{model1}
\begin{tabular}{p{5cm}p{5cm}p{3cm}}
%\multicolumn{3}{c}{\tt Density evolution model }\smallskip\\
\multicolumn{3}{l}{\tt NEUT:;}\smallskip\\
{\tt DN=1+FPR**2;} & {\tt CN=-DN*FX;} & {\tt SNN=SNNEU}\smallskip\\
{\tt NE:Equation;} & {\tt NE=10*sqrt(FPR)+3;} & \smallskip\\
\end{tabular}
\end{equation}
The first line calls the subroutine \tw{NEUT} at each time step. 
The second line defines the diffusion coefficient \tw{DN} as a parabolic
function, the pinch velocity \tw{CN} as a product of this parabola and a
linear function (the minus sign corresponds to the inward pinch). 
The source term \tw{SNN} describes the ionization and recombination
processes making use of the formula \tw{SNNEU}. 
We remind that the same source term could be defined as \tw{SN=SNNEU*NE}.
The only difference between these two representations is in the numerical
implementation, implicit or explicit, respectively. 
If there is no particular requirement then the implicit approximation is
preferable. 
Finally, the third line says that the diffusion equation for \tw{NE}
should be solved with the initial condition defined by the last command.
The boundary condition is not set here explicitly and, therefore, the
value assigned by the initial condition at the point $\rho_B$ will be
used. 
Effectively, it means that the boundary condition \tw{NEB=3} is applied.
Other options for the boundary condition with a prescribed particle flux
or more complicated are shown in the example in Section 5.2.2.

Exactly the same result can be obtained by means of any auxiliary
transport equation \req{Aux}.
$$
\begin{tabular}{p{5cm}p{5cm}p{3cm}}
\multicolumn{3}{l}{\tt NEUT:;}\smallskip\\
{\tt DF1=1+FPR**2;} & {\tt VF1=-DF1*FX;} & {\tt SFF1=SNNEU}\smallskip\\
{\tt F1:Equation;} & {\tt F1=10*sqrt(FPR)+3;} & \smallskip\\
\end{tabular}
$$
Here a diffusion equation is constructed for the dummy variable \tw{F1}
which in all details reproduces the diffusion equation for \tw{NE}
created by the model \req{model1}. 
This example shows how the transport equations for the variables \tw{F1,
F2, F3} can be used for a description of multi-species plasma.

Consider now five different models
\begin{equation}\label{5models}
\begin{tabular}{p{2cm}|p{25mm}|p{25mm}|p{25mm}|p{3cm}}
\tt Model 1&\quad\tt Model 2 &\quad\tt Model 3 &\quad\tt Model 4 
						      &\quad\tt Model 5\\
\tt &\quad\tt TE:As &\quad\tt TE=TEX &\quad\tt TE=TEX &\quad\tt TE=TEX\\
\tt &\quad\tt	    &\quad\tt	     &\quad\tt TE:As  &\quad\tt TE:Eq \\
\end{tabular}
\end{equation}
and assume that a time evolving electron temperature \tw{TEX} is
stored in a data file which will be used together with each of these
models. 
The model 1 is empty.
Nevertheless, nonempty executable code will be created for this model
which makes a number of assignments including those described in
Section 5.3.1.
In particular, the electron temperature will be defined as 
${\tt TE}(\rho)={\tt TEX}(\rho,t_{start})$ where $t_{start}$ is the
initial time of simulation.
The model 2 produces the same effect.
The model 3 makes the assignment ${\tt TE}(\rho,t)={\tt TEX}(\rho,t)$ at
each time so that the complete time evolution of \tw{TEX} will be stored
in \tw{TE}. 
A result of the model 4 is exactly the same as for the model 3.
In the model 5, solving the transport equation is requested,
however, neither transport coefficients nor the power sources are
defined.
For this model, Astra will build a code where the first line \tw{TE=TEX}
is interpreted as the initial condition ${\tt TE}(\rho)={\tt
TEX}(\rho,t_{start})$ and the equation $\prt T_e/\prt t=0$ will be
solved (we assume here that the magnetic field $B_0,$ the plasma density
$n_e$ and the plasma size do not change in time).
Thus, a result of the model 5 will coincide with that of the models 1
and 2. 

One can conclude from this example, that the instruction of the type
\tw{Name:As} is inessential. 
It is really so in most cases, although a formal difference can appear
between model 3 and model 4 if other commands of the type \rf{Eq_com}
are present. 
Namely, with explicit instructions of the type \tw{Name:As}, order of
assignment operators in a code is the same as in a model.
If all these instructions are implicit then the assignment operators
will follow one another as in the sequence 
\tw{NE, TE, TI, CU, F1, F2, F3.}
The logic discussed implies that if the \tw{Type} of evolution is set to
{\tt Assignment} then an assignment expression for this variable is 
expected. 
Although an absence of such an assignment is not considered as an error, 
a warning message is printed for models like the model 2.
Another case when the explicit request \tw{CU:A} is necessary is
discussed below in the example \rf{4models}.

\vskip 2ex
{\footnotesize
The format of equation command line is modified in the most recent
version 5.3 in comparison to the previous versions of the code.
The old form was
$$\mbox{\tt Name:Type\symbol{123}:Expr2\symbol{125}}$$
which can be compared with \rf{Eq_com}.
The qualifier \tw{Type} is here obligatory.
In all other respects the model generated by the old-type equation
command line fully coincides with that generated by the new one.
However, there is a difference in treatment when an equation command
line is absent.
Namely, in old versions, an assignment of time evolution to all main
variables \tw{NE,TE,TI,...} is activated only if the equation
command line with \tw{Type=Assignment} is present.
Coming back to the example \rf{5models} it means that the models 1, 2, 3
and 5 give the same result while only the model 4 gives the time evolving
temperature ${\tt TE}={\tt TEX}(\rho,t).$

}
The next example sets a transport equation for the poloidal flux
$$
\begin{tabular}{p{2.5cm}p{2.5cm}p{2.5cm}}
\multicolumn{3}{r}{\tt Current diffusion equation}	\smallskip\\
{\tt CU:;}	&					\smallskip\\
{\tt CC=CCHR;}	&					\smallskip\\
{\tt DC=DCHR;}	& {\tt HC=HCHR;} &{\tt XC=XCHR;}	\smallskip\\
{\tt CD=CUX;}	& 					\smallskip\\
{\tt CU=1;}	& {\tt MU=.2;}				\smallskip\\
{\tt UEXT=0.2;}	&
\multicolumn{2}{l}{\tt IPL=0.2+0.8*FRAMP(0.,1.);}	\smallskip\\
\end{tabular}
\begin{tabular}{p{7cm}}				\smallskip\\
{\tt !Equation}						\smallskip\\
{\tt !Conductivity \cite{HR}}				\smallskip\\
{\tt !Bootstrap current \cite{Hbs}}			\smallskip\\
{\tt !Pre-calculated driven current}			\smallskip\\
{\tt !Initial distribution}				\smallskip\\
{\tt !Boundary condition}				\smallskip\\
\end{tabular}
$$%\end{equation}
The first line in this model requests the current diffusion equation to
be solved by the code. 
The second line specifies the conductivity \cite{HR} by making use of
the formula \tw{CCHR}.
The third line defines the bootstrap current in a similar way.
The driven current density is supposed to be stored in a data file under
the name \tw{CUX}.
The assignment command \tw{CU=1} is understood as the initial condition.
It sets the flat current distribution normalized to the total current
\tw{IPL}. 
The second command in this line contradicts to the first one.
Therefore, according to the diagram in Section 4.3.2 it is ignored.
The last line includes two contradicting boundary conditions for the
poloidal flux.
According to the priority rule of Section 4.4.2 the second one for the
total plasma current \tw{IPL} is applied and no notice of~~\tw{UEXT=0.2}
is taken here. 

We discuss now different ways of setting the current density profile in
case when the current diffusion equation is switched off.
\begin{equation}\label{4models}
\begin{tabular}{p{2cm}|p{3cm}|p{3cm}|p{4cm}}
\tt Model 1&\qq\tt Model 2 &\qq\tt Model 3 &\qq\tt Model 4\\
\tt CU:AS  &\qq\tt CU=FPR  &\qq\tt CU:As   &\qq\tt CU:A   \\
\tt CU=FPR &\qq\tt         &\qq\tt MU=MUX  &\qq\tt        \\
\end{tabular}
\end{equation}
The first model prescribes parabolic current density and is fully  
equivalent to the second model. 
In other words, the command \tw{CU:AS} can be omitted without 
changing the result.
The same is also applicable to the third model.
This model prescribes a rotational transform profile as stored in
a data file. 
All these three models can also describe time dependent profiles. 
The last model 4 gives the current density profile which at every time
step is defined by the steady state condition \req{icU} as if the
current penetration time (skin time) is much shorter than all the other
characteristic times. 
In this case, the command \tw{CU:A} cannot be omitted.

We remind (see Section~4.3.2) that there is one essential difference in
a way of defining initial current density profile between Model 1 and
Model 3. 
In both cases, the initial setting must be consistent with the total
plasma current.
But, in the first case, the entire current density profile is adjusted,
in the second case, the boundary value of $\mu(\rho_B)$ at the edge
only. 
In other words, the current density in Model 1 is continuous, in Model 2,
it has a surface skin current.

\paragraph{Initial conditions.}
As follows from the examples considered above, assignment statements 
are used for setting initial conditions. 
This implies that an assignment instruction is treated differently
depending on the context. 
Namely, if the left hand side of an instruction is one of main variables
$\{\tw{NE, TE, TI, CU, F1, F2, F3}\}$ and the corresponding equation 
is absent then the assignment is processed as all other assignment
statements with a vector variable on the left hand side.

Otherwise, if the corresponding equation is present then the time
evolution of the quantity is defined by the equation while the
assignment statement is executed only once at the start of simulation.
Moreover, all main variables and few others (see Section~5.3.1) are
defined even if they do not appear in any assignment statement in a
model. 
In the last-mentioned case, data stored in the start file are used.
If no appropriate data are stored there and no assignment is made then
the code results are not relevant. 
In all cases, when an implicit assignment is involved it is applied just
once for the initial time only.

In addition, the initial conditions can be set by a user developed
subroutine. 
In this case, one should beware that an explicit definition in a model
has higher priority than that in a subroutine. 
In other words, if two contradicting settings in a subroutine and in a
model are present then the latter will be used by the code.

\paragraph{Boundary conditions.}
A similar rule is applied to the boundary conditions.
The boundary condition can be set by an explicit assignment command or
by an implicit (Section 4.4) definition.
As in the case of initial conditions, an explicit assignment can be time
dependent while an implicit one is always time independent and assigns a
boundary value at the initial time of the simulation only. 
Once assigned this value does not change afterwards.

This property can be used for more complicated boundary conditions
that cannot be expressed by a simple formula.
Suppose that a subroutine \tw{SOLAY} calculates the electron density at
the edge \tw{NE(NA1)} taking the particle flux at the same point
\tw{QN(NA1)} as an input. 
This can be implemented in a model as the boundary condition
\begin{equation}\label{model2}
\begin{tabular}{p{5cm}p{5cm}p{3cm}}
%\multicolumn{3}{c}{\tt Density evolution model }	\smallskip\\
{\tt NEUT:;} &\multicolumn{2}{l}{\tt SOLAY(QN(NA1),NE(NA1)):}\smallskip\\
{\tt DN=.1;} & {\tt CN=-VP*VRHH;} & {\tt SNN=SNNEU}	\smallskip\\
{\tt NE:;}   & {\tt NE=NEX;}	  & 			\smallskip\\
\end{tabular}
\end{equation}
This model describes a time evolution of the plasma density profile
assuming a constant diffusion and the neoclassical Ware pinch. 
Note that the argument list of a subroutine is not processed by the
Astra compiler. 
Therefore, the abbreviated notations \tw{QNB} and \tw{NEB} can not be
used as arguments of the subroutine.
If an additional instruction like \tw{NEB=NEXB} is present in the model
\rf{model2} then it will override the edge density defined in the
subroutine \tw{SOLAY}.

\paragraph{Shifted boundary conditions.} Additional possibility for
setting boundary conditions is provided by the term in square brackets
in \rf{Eq_com}. 
It enables solving different transport equations on different radial
intervals. 
This option can be useful when the core and periphery transport have
different physical nature and shifting the boundary condition inside
allows to exclude the periphery transport out of consideration.
It can also be useful when the plasma parameters are not known at the
plasma edge and the boundary condition is set at the outermost measured
point. 

According to the syntax of the command \rf{Eq_com} the all three
instructions  
$$
\begin{tabular}{p{4cm}|p{4cm}|p{5cm}}
\tt TI:[0.8*ABC] &\qq\tt TI:[1,0.8] &\qq\tt TI:[2,XFAN(0.8)] \\
\end{tabular}
$$
are equivalent and require solving the heat conductivity equation
for $T_i$ on the interval $0<a<0.8a_B=0.8\tt*ABC.$
Similarly, the instructions
$$
%\begin{equation}\label{3models}
\begin{tabular}{p{4cm}|p{4cm}|p{5cm}}
\tt TI:[2,0.8] &\qq\tt TI:[AFX(0.8)] &\qq\tt TI:[1,AFX(0.8)/ABC] \\
\end{tabular}
%\end{equation}
$$
require solving the same equation on the interval $0<\rho<0.8\rho_B.$

The next model
\begin{equation}\label{model3}
\begin{tabular}{p{6cm}p{6cm}}
\multicolumn{2}{l}{\tt IFSPPL(CAR1,CAR2,CAR3,CAR4,CAR5,CAR6):}
\smallskip\\
{\tt XI=HNCHI+CAR5;} & {\tt PI=PEICL+PIX}
\smallskip\\
{\tt TI:Eq[1,0.9];}	& {\tt TI=TIX}
\smallskip\\
\end{tabular}
\end{equation}
produces a transport code solving the heat conductivity equation for
$T_i$ on the interval $0\leq a\leq 0.9a_B.$ 
Here the ion heat conductivity $\chi_i=\tt XI$ is defined as a sum of
neoclassical and anomalous terms where the latter (along with many other
quantities) is calculated by the subroutine \tws{IFSPPL}. 
The heating source \twd{PI} includes the electron-ion heat exchange
\twd{PEICL} and the auxiliary heating \twd{PIX} defined in a data file.

The command \twd{TI=TIX} bears here a double meaning.
Firstly, it defines the ion temperature $T_i(a,t_0)=T_{i,exp}(a,t_0)$ at
the initial time $t_0$ on the whole radial interval $0\leq a\leq a_B.$ 
The subsequent time evolution of $T_i(a,t>t_0)$ on the partial interval
$0<a<0.9a_B$ is defined by the heat conductivity equation. 
Thus, on the segment $0\leq a\leq 0.9a_B,$ the command \twd{TI=TIX}
provides the initial and the boundary condition\footnote
{Like in the case of non-shifted boundary, this assignment is made
only once at the start time $t=t_0.$ 
Moreover, it applies only when no explicit boundary setting is provided.}
for the corresponding transport equation. 
Secondly, it defines the time evolution of the ion temperature 
$T_i(a,t)$ on the interval $0.9a_B<a\leq a_B$ for $t>t_0$ that is not
covered by the transport equation.
Therefore, the command \twd{TI:[1,0.9]} in the last line of the model
\rf{model3} can be thought of as if it combines the definition of
\twd{Type=Equation} for the interval $0<a<0.9a_B$ and of
\twd{Type=Assignment} outside it. 

Similarly to the case of non-shifted boundary conditions, the boundary
value $T_i(0.9a_b)$ implicitly imposed by the command \twd{TI=\dots} is
time independent. 
In order to set a time dependent boundary condition one has either to
use an explicit assignment statement in any of the forms
$$
\begin{tabular}{p{4cm}|p{4cm}|p{5cm}}
\tt TIB=... &\qq\tt QIB=... &\qq\tt QITB=\dots \\
\end{tabular}
$$
or to implement an implicit boundary setting with an external subroutine
similar to that of the example \rf{model2}.

It is significant to note that, in the particular case of the shifted
boundary conditions, the abbreviation \tw{TIB} is interpreted
distinctively depending of the context.
Namely, whenever \twd{TIB} appears on the right hand side of any
expression, it is replaced by the Astra code compiler with \twd{TI(NA1)}
which is equivalent to $T_i(a_B).$ 
However, if (and only if) any of the notations \twd{NEB, QNB, QNNB,
TEB} and so on\footnote{The complete list is given in Table~4.8.} 
stands on the left hand side of an assignment command then this command
is interpreted as the boundary condition assignment and it is attributed
to the boundary point declared by a corresponding equation instruction
of the type \rf{Eq_com}. 

For instance, if the command \twd{TIB=TIXB} is used in the model
\rf{model3} it will result in the Fortran statement
\twd{TI(NA1I)=TIX(NA1)}, where the array index \twd{NA1I} points to the
grid node nearest to $0.9\,a_B$ while the index \twd{NA1} to the edge
grid point $a_B=\tt ABC.$
Most probably, it is not the result expected by the user.
The correct assignment can be implemented as
\twd{TIB=TIX(0.9*ABC).}

\subsubsection{Radial output}
First of all we note that there is no default output in Astra
and therefore quantities of interest must be explicitly specified in the
model definition.
Request for a radial output has the form
\begin{equation}\label{Radout}
\mbox{\tt 
Name\bs\symbol{123}Expression\symbol{125}%
\symbol{123}\bs\bs Xdata\symbol{125}%
%\symbol{123}\%Xdata\symbol{125}%
\symbol{123}\bs Scale\symbol{125}%
%\symbol{123}[Expr1,Expr2]@\#\%Tag\symbol{125}
}\qquad
\end{equation}
Here \twd{Name} specifies a legend (a string with a length up to 4
characters) associated with the radial profile given by any regular
arithmetic \twd{Expression}. 
Every part of this command shown in braces is optional and can be
omitted, the group preceded by \twd{\bs\bs} can either be omitted or 
repeated several times. 

Up to 96 commands of the type \rf{Radout} can be present in a model
and, consequently, up to 96 radial profiles can be plotted (or stored)
in a simulation.
Each radial profile defined by the corresponding \twd{Expression}
together with its \twd{Name} and \twd{Scale} will appear in a graphic
window when the code starts execution.
The code provides a few different graphic modes which allow combination
of several curves in one box and several boxes in one window.
The curves are sequentially placed in boxes in the same order as the
curve descriptions \rf{Radout} appear in the model. 
If a number of curves is larger than the number of boxes then the next
curve is again put in the first box, then in the second and so on until
all boxes in the current window acquire the maximal for this mode number
of curves. 
Then the next (hidden) window is filled and the process continues until
the list of the curves submitted to output is exhausted.

In other words, one can say that each radial profile defined with
\rf{Radout} receives its sequence number which is an ordinal number of
the output command as they appear in the model.
The curves are placed in graphic windows according to their sequence
numbers.
Because the position of the curves on the display it is essential for
the convenience of viewing results, it is important to pay attention to
the sequence of the output commands.
Nevertheless, the user can change the sequence numbers and other curve
attributes interactively during the code execution. 
More detailed description of this feature will be given in Section~5.5.3.

The parameter \twd{Scale} defines the scale for plotting the curve given
by \twd{Expression}. 
If \twd{Scale} is zero or absent then the scale of the output profile is
automatically selected by the code. 
If \twd{Scale} is given as a positive number then this number specifies
the actual scale for the corresponding radial profile.
Negative integer value of \twd{Scale} allows setting the same automatic
scale for all profiles in the group with the same negative scale value.
It is useful for representation of several different curves of the
same meaning, e.g. the input powers of different heating methods.

There is one more peculiarity in setting automatic group format.
Normally auto-scale is selected based on the maximum absolute value in
the entire group.
But the user can enforce the code to make automatic scale basing upon a
single specified profile in the group.
It happens if the absolute value of the negative scale coincides
with the sequence number of the output.
We illustrate the discussed rules for radial output with an example:
\vskip 1ex\noindent
{\bf Example:}\smallskip\\ %{\rm(1)}\\
{\tt\hspace*{1mm}
Te\bs TE\bs -9;~~~~~~~Ti\bs TI\bs -9;~~~~~~~j\bs CU\bs-2;~~~~~~~~
		Utor\bs UPL\bs-1;\smallskip\\ \hspace*{1mm}
ne\bs NE\bs 10;~~~~~~~\bs;~~~~~~~~~~~~~~jOhm\bs CUOHM\bs-2;~~
		iota\bs MU;\smallskip\\ \hspace*{1mm}
Texp\bs TEX\bs-9;~~~~Tixp\bs TIX\bs -9;~~~~jBS\bs CUBS\bs-2;~~~~
		Ulng\bs ULON\bs -1;\smallskip\\ \hspace*{1mm}
nexp\bs NEX\bs 10;~~~~\bs;~~~~~~~~~~~~~~sigm\bs CC;~~~~~~~~~q\bs 1/MU;
}\vskip 1ex

%\begin{equation}\label{Rexa}
%\mbox{\tt\begin{tabular}{p{3cm}p{3cm}p{3cm}p{3cm}}
%Te\bs TE\bs -9; & Ti\bs TI\bs -9;  & j\bs CU;	  & Utor\bs UPL; \\
%ne\bs NE\bs 10;  & \bs; jOh\bs CUOHM; &   & iota\bs MU; \\
%Texp\bs TEX\bs-9;&Tixp\bs TIX\bs -9;&jBS\bs CUBS;&Ulng\bs ULON\bs-1;\\ 
%nexp\bs NEX\bs 10;& &\bs; sigm\bs CC;         & q\bs 1/MU;\\
%\end{tabular}}
%\end{equation}

\noindent
Assume that this set of data is presented in the default radial output
mode which allows up to 16 curves in a window. 
The window includes 8 boxes placed in two rows and four columns with
two curves in a box as schematically shown in this sketch.
\bigskip\\
\cl{\tt
\begin{tabular}{| l | l | l | l |}				\hline
\bf\large 1 &\bf\large 2 &\bf\large 3 &\bf\large 4 \\
1 \& 9		&2 \& 10	&3 \& 11	&4 \& 12	\\ \hline
\bf\large 5 &\bf\large 6 &\bf\large 7 &\bf\large 8 \\
5 \& 13		&6 \& 14	&7 \& 15	&8 \& 16	\\ \hline
\end{tabular}
}
\smallskip\\
Hear the large digits in the left upper corner of each box show the box
numbers and the curve sequence number are shown in ordinary type.
Then the first eight profiles will be shown in red and the next eight in
blue so that each box will contain one red and one blue curve.
Consequently, the first and the third lines in the example enter in the
first row of boxes while the second and the forth lines occupy the second
row. 
The empty instructions \twd{\bs;} are used in this example in order to
skip the sixth box and put all quantities related to the current density
one over or under another for a more convenient comparison.

Thus the first box contains calculated and experimentally
measured electron temperatures, \twd{TE} and \twd{TEX}, respectively. 
The former is labeled on the plot as \twd{Te} and the latter as 
\twd{Texp}. 
Both will have the same scale defined by the data \twd{TEX} (sequence
number~9).
Similarly, the next box contains two ion temperatures with the same
scale. 
The third box includes the total current density \twd{CU} and the
bootstrap current density \twd{CUBS} each with a common scale.
The last box in the first row encloses the toroidal \twd{UPL} and the
longitudinal \twd{ULON} loop voltages in the same scale defined by the
largest of the two.
The fifth box shows two density profiles with the
prescribed scale. 
The next box (No.~6) is empty. 
The seventh box shows the Ohmic current density and conductivity
profiles (which must be proportional in the steady state).
Finally, the last in this mode, eighth box presents the rotational
transform \twd{MU} and the safety factor.

\paragraph{Presentation of experimental data.}
So far we discussed the output of radial profiles presented as curves.
In the example above, even so called ``experimental'' data as \twd{TEX}
or \twd{NEX} are smoothed and transferred to the fine transport grid
with a typical size of 100 nodes. 
Transferring input experimental data to a radial grid of the 
transport code is a problem which has no universal solution. 
Although the Astra code provides a way of doing this transfer
it could be unsatisfactory in some cases, for instance, it destroys
a significant information while smoothing data or, in opposite, add 
artificial features while extrapolation.
In order to provide a way of visual control of data treatment
the Astra code enables displaying the original input data as the
radial plots. 

Such a plot can be requested by a command line of the format \rf{Radout}
making use of the control group \twd{\symbol{34}\bs\bs\symbol{34}}.
A quantity \twd{Xdata} will be shown on a plot as it is stored in the
input file (Section~5.4). 
To be interpreted properly, the control word \twd{Xdata} must be one of
the array names given in Table~5.4.
For instance, the instruction \\
\hspace*{5cm}{\tt Tex\bs TEX\bs\bs TEX\bs 10}\\
plots the electron temperature \twd{TEX}, firstly, as a smoothed curve,
secondly, by dots as the original data are stored in the input file.
The curve and the dots will appear in the same box with the same
scale 10 keV. 
One can plot the dots only \\
\hspace*{5cm}{\tt Tex\bs\bs TEX\qquad\qquad\qquad}\\
or combine several measurements in one plot\\
\hspace*{5cm}{\tt Tex\bs TEX\bs\bs TEX\bs\bs CAR1X}\\
The latter command plots \twd{TEX} as a curve together with two sets of
data stored in the input file under the names \twd{TEX} (e.g. Thomson
scattering) and \twd{CAR1X} (e.g. ECE data). 
The data for $\tt TEX$ and $\tt CAR1X$ will be shown by different
symbols. 

%\paragraph{Radial scale control.}
%Sometimes it is needed to inspect a fine structure of radial dependence 
%that cannot be clearly resolved in the standard radial scale.
%To do this one can extend a part of the radial axis to the whole box.
%It can be done by adding an instruction

\subsubsection{Time output}
Radial and time output commands have a similar format.
The difference is that the backslash symbol 
\twd{\symbol{34}\bs\symbol{34}} is replaced by the underscore
\twd{\symbol{34}\us\symbol{34}}. 
This format reads
\begin{equation}\label{Timout}
\mbox{\tt\symbol{123}Name\symbol{125}\us\symbol{123}Expression\symbol{125}\symbol{123}\us Scale\symbol{125}}\qquad
\end{equation}
Similar to the radial output, any allowed arithmetic expression can be
used for output. 
However, a natural restriction is applied that only scalar (radially
independent) expressions may be used in the time output instructions. 
The user can prescribe an output scale, ask for a common scale for several
plots or take advantage of automatic scale adjustment. 
Up to 96 time dependences can be processed by the current version of the
code. 

\subsubsection{Advanced options}

\paragraph{User defined variables.}

Variables used in a transport model basically belong to the Astra
internal variable list which includes (a) a set of physically meaningful
predefined variables and (b) a set of holder scalar and vector variables
without specific meaning. 
It is not considered as an error if the user employs his own scalar
variables in a model. 
Though there is no support for this variables and they are not defined 
in most of Astra modules except for if explicitly passed to a
corresponding subroutine.

For instance, one can define a time dependent scalar quantity\\
\hspace*{15mm}
\twd{STAIRS=FJUMP(0.1)+FJUMP(0.2)+FJUMP(0.3)+FJUMP(0.4)+FJUMP(0.5)}\\ 
and later use it as a part of another assignment statement or as a
parameter of a subroutine, however, the quantity \twd{STAIRS} cannot be
plotted with the standard time output instruction~\rf{Timout} because it
is not available in the module concerned with the output.
In order to overcome this limitation the user has two options.
The standard and recommended way is to replace the name \twd{STAIR} with
one of the C-parameters (Table~4.17).
Another option is to describe the variable \twd{STAIRS} in a common
block inside the file\footnote{See footnote on the page \pageref{AWD}.}
\\ 
\hspace*{3cm}
\twd{AWD/tmp/declar.usr} \\
This file is included by the Fortran statement \twd{INCLUDE} in all
model dependent modules of the Astra kernel.

\paragraph{Verbatim reproduction.}

Processing a model file the Astra compiler transforms it into a Fortran
source code according to some rules. 
However, sometimes these rules are too restrictive because they are
invented exclusively for dealing with Astra variables.
There is a possibility to override these rules, switch off all
transformations and require a verbatim record into a source code. 
It can be done by quoting the appropriate string. 
For instance, the expression 
\\
\hspace*{4cm}{\tt rhNA\us\symbol{34}RHO(NA)\symbol{34}}
\\
yields a value of the \tws{NA}-th element of the radial grid array
\twd{RHO,} i.e. the position of the grid point adjacent to the edge
one $\rho_B=\tt RHO(NA1).$
Similarly, the instruction 
\\
\hspace*{4cm}\tw{wrk3\bs\symbol{34}WORK(j,3)\symbol{34}}\\
allows to plot the data stored in the 3rd column of the two-dimensional
array \twd{WORK}.

It is not out of place to stress here that a care should be taken when
using this feature.
The user should beware that the Astra compiler can put the quoted
string in unexpected place of the code or even being placed in a
proper position the string can get in conflict with Fortran rules.
Therefore, it is recommended to check if the effect of such a
nonstandard operation provides the desired result.

In addition we add a remark concerning the two working arrays
\twd{WORK} and \twd{WORK1}. 
Both arrays are described as real two-dimensional arrays
and placed in a common block\footnote
{In the standard installation $\twd{NRD}=501.$} 
\\
\hspace*{1.5cm}\tw{COMMON~}
\tw{/WORKAR/~WORK1(NRD,2*NRD+7),~WORK(NRD,2*NRD)}\\
in the file \twd{AWD/for/status.inc} 
so that they can be used in every module including this file.
The arrays are not being used by any routine of the Astra kernel.
It is supposed that the first array \twd{WORK1} is used as a working
area for plug-in subroutines without any usage outside these subroutines.
Another array \twd{WORK} may be used either as a working area or as a
tool for data exchange between user's subroutine and other parts of
the code. 
In the second case, it remains at user's discretion to trace possible 
conflicts between different user developed modules.

\paragraph{Post-compiler editing.}
A possibility to interfere in the pre-defined sequence of Astra
operations is given by the command
\\
\tw{> Astra -t}
\\
This suspends the code execution after the second step (see the flow
diagram at page \pageref{flow}).
At this point, the user can introduce any changes into the changeable
part of the code.
This part consists of intermediate Fortran source files located in the
directory \twd{AWD/tmp/}. 
When the appropriate adjustments are performed the code execution will
be resumed with account for all recent modifications.

\subsubsection{Several examples}
In this section we consider several examples of building different
models in the Astra system.
We assume that all models discussed below are run with an appropriate
input file that provides all data required by a model. 
Therefore, some properties of data setting are used already in this
section although they are first introduced in Section~5.3.

We start with a discussion of simple plasma characteristics which can be
calculated on the basis of measured plasma density and temperature
profiles.  
\vskip 1ex\noindent
%{\bf Example 1:}\smallskip\\ %{\rm(1)}\\
{\tt\hspace*{1mm}
Wth\bs WTOT;~~~~~LTi\bs LTI;~~~~~~~Heff\bs HEEFF;~~~~Vpol\bs VPSWW;
\smallskip\\ \hspace*{1mm}
tauE\bs TAUE;~~~~LTic\bs RLTCR;~~~~Hiff\bs XIEFF;~~~~rhos\bs RLS;
\medskip\\ 
}
As a result of this model the following quantities will be plotted one
under another (in the 1st mode, their placement on a screen is the same
as in the printout above): 
the plasma thermal energy content (possible contribution from fast
particles is not included) and the energy confinement time; 
inverse normalized length of the ion temperature mapped to the
mid-plane, $R_0/L_{Ti}=R_0|\prt T_i/\prt a|/T_i$;
and a critical value for the same quantity with respect to the ion
temperature gradient instability;
effective (i.e. experimentally measured) heat conductivities for the
electron and ion thermal components;
and, finally, the neo-classical poloidal rotation velocity and 
$\rho_s=\sqrt{T_e/m_i}.$
All they will be plotted as radial functions.

This code will also work if some quantities needed for calculations are
not defined, however, then the results can be irrelevant.
For instance, if the heating power is not defined then the heat
conductivities and the confinement time will be negative.
In order to get appropriate values for these quantities one has to define
the input heating power in any allowed way, e.g. as 
\twd{PE=POH;~PI=PIX}.

Another essential remark is that the radial functions plotted by this
model are time independent even if the correspondent data file contains
time dependent data. 
This happens in line with the syntax rules described in Section~5.2.2.
Indeed, without the explicit definition the assignment as \twd{TE=TEX}
and similar are fulfilled only once at $t=t_0,$ consequently, all
quantities calculated on the basis of $n_e,\,T_e$ and so on are time
independent. 
One can view the difference running the model
\vskip 1ex\noindent
{\tt\hspace*{1mm}
Tex\bs TEX\bs\bs TEX\bs -9;~~~~2\bs;~~~3\bs;~~~4\bs;
	    ~~~~~~~~~~~~~5\bs;~~~6\bs;~~~7\bs;~~~8\bs;
\smallskip\\ \hspace*{1mm}
Te\bs TE\bs -9;
\medskip\\}
Two graphs will show the quantities \twd{TE} (in blue) and \twd{TEX}
(in red) as smoothed curves in the same scale in the upper left box. 
Additionally, the original data for \twd{TEX} will be shown with red
crosses.
The curve for \twd{TEX} will evolve as it is written in the data file,
while \twd{TE} will coincide with \twd{TEX} at the first instant only.

In order to remove all differences, one has to include in the model the
explicit definitions of main plasma parameters.
\vskip 1ex\noindent
{\tt\hspace*{1mm}
NE=NEX;~~~TE=TEX;~~~TI=TIX;~~~~~~~NEQUIL=41;
\medskip\\}
Here the command \twd{NEQUIL=41} requires that the Astra built-in 
3-moment equilibrium solver is employed for the equilibrium 
reconstruction.\footnote
{We remind that, if not defined explicitly, the
current density distribution is calculated from the steady state
condition \req{Uplss}.}

As an example of the time output consider
\vskip 1ex\noindent
%{\bf Example 2:}\smallskip\\ %{\rm(1)}\\
{\tt\hspace*{1mm}
ABC\us ABC;~~~~~~kapp\us ELONG;~~~~~~delt\us TRIAN;~~~~~~shif\us SHIFT;
\smallskip\\ \hspace*{1mm}
Tex0\us TEXC;~~~~Tix0\us TIXC;~~~~~~~Ipl\us IPL;~~~~~~~~~Btor\us BTOR;
\smallskip\\ \hspace*{1mm}
Wth\us WTOTB;~~~~pk1\us NEC/NEAVB;~~~pk2\us NEC/NECHC;
				~~H-89\us TAUEB/(TITER+1.e-6);
\smallskip\\ \hspace*{1mm}
q0\us 1/MUC\us 5;~~~qa\us 1/MUB\us 5;~~~~~~qmin\us QMINB\us 5;
		~~~qmin\us 1/FRMAX(MU)\us 5;
\smallskip\\ \hspace*{1mm}
Xmin\us XQMINB\us 1;~~~~q2a\us AFVAL(MU,.5)/ABC\us 1;
		~~~~~~q3a\us AFVAL(MU,.333)/ABC\us 1;
\medskip\\}
The first line shows the geometrical parameters of the outermost
magnetic surface: the minor radius \twd{ABC}, elongation \twd{ELONG},
triangularity \twd{TRIAN} and shift with respect to the major radius
\twd{RTOR}.
The second line gives the central values of electron and ion
temperatures then the plasma current and magnetic field. 
After that follow: the total plasma energy content, the density peaking
factor calculated as the central value divided by the volume averaged
and by the chord averaged densities, enhancement factor with respect to
the ITER-89 scaling. 
The next two lines show different quantities characterizing the safety
factor $q$ distribution. 
The first of them yields the central $q(0)$ and the edge $q(a_B)$ values,
then the minimum $q_{min}$ value calculated in two different ways. 
Finally, the last line gives the radial positions (in the normalized 
coordinate $x=a/a_B$) of the points where $q(x)$ reaches its minimum 
value, where $q(x)=2$ and $q(x)=3.$

Unlike the radial output most of the time dependencies above are plotted
as they are stored in the data file, i.e. if a time dependence of any
Astra simple variable (e.g. \twd{ABC}, \twd{IPL}, \twd{BTOR}, \dots) is
written in the data file then it will be shown as time varying also in
the Astra output. 
This can be not the case for the derived quantities.
For instance, \twd{NEC} and \twd{NEAVB} are calculated making use of
\twd{NE}, therefore, for those quantities, the same rules are applicable
as for \twd{NE}.

The next example shows a model which reads in the experimental data for
$n_e,T_e,T_i$ and calculates the current density evolution assuming
the neoclassical conductivity and bootstrap current.
\vskip 1ex\noindent
{\tt\hspace*{1mm}
NE=NEX;~~~~~TE=TEX;~~~~~TI=TIX;~~! measured plasma parameters
\\ \hspace*{1mm}
AMAIN=2;~~~~ZMAIN=1;~~~~~~~~~~~~~! mass \& charge of the main ion species
\\ \hspace*{1mm}
AIM1=12;~~~~ZIM1=6;~~~~~~~~~~~~~~! mass \& charge of the main impurity
\\ \hspace*{1mm}
NIZ1=NE*(ZEF-1)/(ZIM1-1)/ZIM1;~~~! density of the impurity
\\ \hspace*{1mm}
NI=NE*(ZIM1-ZEF+1)/ZIM1;~~~~~~~~~! total density of all plasma ions
\\ \hspace*{1mm}
NDEUT=NI-NIZ1;~~~~~~~~~~~~~~~~~~~! density of deuterium plasma ions
\\ \hspace*{1mm}
NEQUIL=41;~~~~~~~~~~~~~~~~~~~~~~~! 3-moment equilibrium solver
\\ \hspace*{1mm}
HC=HCKIM;~~~DC=DCKIM;~~~XC=XCKIM;! Bootstrap due to Kim
\cite{Kim} \\ \hspace*{1mm}
CU:EQ;~~~~~~MU=.33*FPR+.33;~~~~~~! current density Eq \& initial
condition 
\\ \hspace*{1mm}
CC=CNHR;~~~~CD=CUBM;~~~~~~~~~~~~~! conductivity \cite{HR}~ \& NB driven
current 
\\ \hspace*{1mm}
MIXINT(CV1,CF1):;~~~~~~~~~~~~~~~~! Kadomtsev reconnection
\\ \hspace*{1mm}
NBI:;~~~~~~~~~NEUT:;~~~~~~~~~~~~~! NB injection and cold neutrals
\smallskip\\ \hspace*{1mm}
~~~~~~~~~~~~~~~ Radial output
\smallskip\\ \hspace*{1mm}
Tex\bs TEX\bs\bs TEX;~~~~~q\bs 1./MU\bs 10;
			~jtot\bs CU\bs 2;~~~~~Upl\bs UPL;
\\ \hspace*{1mm}
nex\bs NEX\bs\bs NEX;~~~~~cneo\bs CC;
			~~~~jNB\bs CUBM\bs 2;~~~~shir\bs SHEAR;
\\ \hspace*{1mm}
Tix\bs TIX\bs\bs TIX;~~~~~iota\bs MU;
			~~~~jOH\bs CUOHM\bs 2;~~~ULON\bs ULON;
\\ \hspace*{1mm}
Zeff\bs ZEFX\bs\bs ZEFX;~~cSp\bs CCSP;
		~~~jBS\bs CUBS\bs 2;~~~~psiN\bs (FP-FPC)/(FPB-FPC);
\smallskip\\ \hspace*{1mm}
~~~~~~~~~~~~~~~~Time output
\smallskip\\ \hspace*{1mm}
Upl\us UPLB\us 2;~~~betj\us BETAJB;~~~bett\us BETTB;
		~~~~li\us LINTB*RTOR/(RTOR+SHIFT)\us 1;
\\ \hspace*{1mm}
Ipl\us IPL\us 2;~~~~Ibs\us IBSB\us 2;
		~~~Iohm\us IOHMB\us 2;~~~Inb\us IBMB\us 2;
\medskip\\}
This model can be used for simulating NBI assisted current ramp-up phase
in a tokamak.
Therefore, most of output quantities are related to the current density
distribution. 
In particular, the radial output shows the neoclassical and Spitzer
conductivities (6th box), the third column (3rd and 7th boxes) includes
different contributions to the current density, the forth box exhibits a 
difference between the toroidal and longitudinal loop voltages, the
shear and the normalized poloidal flux are displayed in the 8th box.
The time output presents the edge loop (toroidal) voltage, poloidal and
toroidal betas, internal inductance per unit length, the total plasma
current as comprised of the bootstrap, Ohmic and beam driven currents.

The next example illustrates simulation of a plasma density build-up.
\smallskip\\
%\vskip 1ex\noindent
{\tt\hspace*{1mm}
~~~~~~~~Equilibrium, electron, ion temperatures and poloidal flux
\\ \hspace*{1mm}
NEQUIL=41;~~~~~~~TE=2*FPR**2+.1;~~~~~~TI=TE/2;~~~~~~CC=CCSP;
\smallskip\\ \hspace*{1mm}
~~~~~~~~Particle source (wall neutrals + pellets)
\\ \hspace*{1mm}
NEUT:;~~~~~~~~~~~SNN=SNNEU;~~~~~~~~~~~~~! Gas puffing
\\ \hspace*{1mm}
CV1=anint(.1*sin(100*GP2*TIME)+.405);~~~! Square wave-form oscillator
\\ \hspace*{1mm}
SN=CF3*30000*CV1*GAUSS(.7,.1);~~~~~~~~~~! Pellet evaporation profile
\smallskip\\ \hspace*{1mm}
~~~~~~~~Density eqn with initial and boundary conditions
\\ \hspace*{1mm}
NE:EQ;~~~~~~~~~~~NE=10*FPR;~~~~~~~~~QNNB=CF2*10000;
\\ \hspace*{1mm}
CAR1=CF1*(1+3*FX**2);~~~~~~~~~~~~~~~~~~~! Prescribed function
\\ \hspace*{1mm}
DN=CAR1*(1-XSTEP(.85))+.5*XSTEP(.85);~~~! Diffusion coefficient
\\ \hspace*{1mm}
CN=-VP*VRHH;~~~~~~~~~~~~~~~~~~~~~~~~~~~~! Neoclassical pinch
\\ \hspace*{1mm}
TAUMAX=.0001;~~~~~~~~~~~~~~~~~~~~~~~~~~~! Time step limitation: < 1 ms
\smallskip\\ \hspace*{1mm}
~~~~~~~~~Auxiliary quantities
\\ \hspace*{1mm}
CV16=VINT(SNTOTB);~~~~~~~~~~~~~~~~~~~~~~! The same as QNTOTB
\\ \hspace*{1mm}
CV15=TIMINT(CV16);~~~~~~~~~~~~~~~~~~~~~~! Total particle source integral
\\ \hspace*{1mm}
CV14=NEAVB;~~~~~~~~~~~~~~~~~~~~~~~~~~~~~! Volume average density <ne>
\\ \hspace*{1mm}
CV13=-TIMDER(CV14);~~~~~~~~~~~~~~~~~~~~~! d<ne>/dt
\\ \hspace*{1mm}
CV2=TIMINT(CV1);~~~~~~~~~~~~~~~~~~~~~~~~! Total evaporation time
\smallskip\\ \hspace*{1mm}
~~~~~~~~~Radial output
\\ \hspace*{1mm}
D\bs DN;~~~~~~~~gn\bs GN;~~~~~~~flux\bs QN;~~~~~~~~etai\bs LNE/LTI\bs
								100;
\\ \hspace*{1mm}
ne\bs NE;~~~~~~~NN\bs NN\bs 1;~~~~~nue*\bs NUES\bs 1;~~~~Spuf\bs SNN*NE;
\\ \hspace*{1mm}
Vin\bs -CN;~~~~~V/D\bs -CN/DN;~~I(S)\bs QNTOT;~~~~~Lne\bs LNE;
\\ \hspace*{1mm}
Stot\bs SNTOT;~~TN\bs TN;~~~~~~~Ware\bs VP*VRHH;~~~Spel\bs SN;
\smallskip\\ \hspace*{1mm}
~~~~~~~~Time output
\\ \hspace*{1mm}
n(0)\us NEC;~~~~~~~tauP\us TAUPB\us .1;~~~~~~nlin\us NECHC;
						~~~~~~~~Gin\us CV16;
\\ \hspace*{1mm}
n(a)\us NEB;~~~~~~~tau\us CV14/CV13\us .1;~~~<ne>\us CV14;
						~~~~~~~~~Gout\us QNB;
\\ \hspace*{1mm}
<n>G\us 2.7*IPL/ABC**2;~~~NNCL\us NNCL;~~~ntot\us VOLUME*CV14;~~~NNWM\us NNWM;
\\ \hspace*{1mm}
Sinp\us CV15;~~~~~~~~~~~~~albe\us ALBPL;~~1or0\us CV1;~~~~~~~~~~~duty\us CV2/TIME;
\medskip\\}
Only the density diffusion equation is solved in this model.
The pinch is assumed to be neoclassical while the particle diffusion is
taken as a parabolic function inside $\rho<0.85\rho_B$ (plasma core) and
as a constant $\rm 0.5~m^2/s$ outside this region (pedestal zone).
The boundary condition for the density is given by the requirement that
the total particle flux through the plasma boundary is
$Q_n=V<n_e>/\tau_P=\tw{QNNB*}n_e(a).$
For ASDEX-Upgrade parameters this gives 
$Q_n\approx 3\times 10^{22}\rm\;s^{-1}.$

The particle source consists of two comparable parts: gas puffing and
pellet injection. 
The former is described in Section~4.9.4. 
The latter is assumed to be periodic ($1/T=100$Hz) with evaporation time
$\Delta t\approx 1$ms. 
The pellet ablation profile is represented by a Gaussian function with
the prescribed width $\Delta_\rho=0.1\rho_B$ and the maximum position 
$\rho_0=0.7\rho_B$. 
The amplitude of this source ($3\times 10^{23}\rm\;s^{-1}$ in our example) can be evaluated as 
$Q_nT/\Delta t.$
Note that no interaction between the wall and pellet neutrals is taken
into account in this simplified model.
%Here \twd{CF1, CF2, CF3} are fit coefficients all of order 1.

The radial output in this model gives only the quantities related to the 
particle behavior. 
It shows the diffusion coefficient \twd{DN}, the pinch velocity \twd{CN}
(positive sign means outward flux), the total particles source
\twd{SNTOT} and the partial sources due to pellets \twd{SN} and cold
neutrals \twd{SNN*NE}.
The quantity $\nu_e^*=\tt NUES$ which characterize collisions is
shown together with the Ware pinch velocity \twd{VP*VRHH}.
Also of interest could be a comparison of the two quantities \twd{QN},
collisional particle flux, and \twd{QNTOT}, ionization particle source.
In the steady state, this two quantities must coincide.
In this simulation, they are always strongly different.

The time output shows central, edge, volume and line average densities.
The particle confinement time is calculated in two different ways.
The Greenwald limiting density is shown under the name \twd{<n>G}.
The two components of incoming cold neutrals are characterized by the
parameters \twd{NNCL} and \twd{NNWM}.

We conclude this section with a simulation model for electron and ion 
power balance based on the IFS/PPPL transport model \cite{ifspppl}.
For more convenient discussion we split the model in two pieces.
\\
{\tt\hspace*{1mm}
!=============== ITG instability based transport model (part 1) =========
\\ \hspace*{1mm}
NEQUIL=41;~~~~~~~~~~~~!~~~3-moment solver for the Grad-Shafranov equation
\\ \hspace*{1mm}
!-------------------- Electron and ion densities ------------------------
\\ \hspace*{1mm}
NE:AS;~~~~~~~~NE=NEX;~~~~~ZEF=ZEFX;
\\ \hspace*{1mm}
AMAIN=AMJ;~~~~ZMAIN=ZMJ;~~~~~~~~~~~!~mass \& charge of main ion species
\\ \hspace*{1mm}
AIM1=12.;~~~~~ZIM1=6.;~~~~~~~~~~~~~!~mass \& charge of 1st impurity
\\ \hspace*{1mm}
ZEF1=ZIM1*(ZEF-1.)/(ZIM1-1.);~~~~~~!~1st impurity contribution to Zeff
\\ \hspace*{1mm}
NIZ1=NE*(ZEF-1.)/(ZIM1-1.)/ZIM1;~~~!~density of 1st impurity
\\ \hspace*{1mm}
NI=NE*(ZIM1-ZEF+1.)/ZIM1;~~~~~~~~~~!~density of all ions
\\ \hspace*{1mm}
NDEUT=NI-NIZ1-NIBM;~~~~~~~~~~~~~~~~!~density of deuterium plasma ions
\\ \hspace*{1mm}
!-------------------- Auxiliary heating ---------------------------------
\\ \hspace*{1mm}
NBI:0.01;~~~~~~~~~~~~~~~~~~~~~~~~~~!~Neutral beam injection
\\ \hspace*{1mm}
!-------------------- Plasma rotation -----------------------------------
\\ \hspace*{1mm}
VTOR=VTORX;~~~~~~VPOL=VPSWW;~~~~~~~~~~~!~VTOR from exp., VPOL from neocl.
\\ \hspace*{1mm}
ER=BTOR*(FRS*MU*VTOR/RTOR+VDIA-VPOL);~~!~Radial electric field 
\\ \hspace*{1mm}
!-------------------- ISF/PPPL transport model --------------------------
\\ \hspace*{1mm}
!~~~~~~~~~~~~~~~~~~~~~~~~~Output of IFSPPL:
\\ \hspace*{1mm}
!~~(1) R/L\us Tcrit for ITG mode,~~(2) R/L\us Tcrit for C branch,
						~(3) not used,
\\ \hspace*{1mm}
!~~(4) Mode increment,
~~~~~~~~~(5) CAR24=chi\us i,~~~~~~~~~~~(6) CAR23=chi\us e
\\ \hspace*{1mm}
IFSPPL(CAR19,CAR20,CAR21,CAR22,CAR24,CAR23):;
\\ \hspace*{1mm}
SMEARR(CV1,CAR23,CAR25):;~~CAR17=FTAV(CAR25,CV2);~~!~Space and time
\\ \hspace*{1mm}
SMEARR(CV1,CAR24,CAR26):;~~CAR18=FTAV(CAR26,CV2);~~!~~~averaging
\\ \hspace*{1mm}
HE=CAR17+HNGSE;~~~~~~~~~~~~XI=CAR18+FOWC*HNCHI;~~~~!~Heat conductivity
\\ \hspace*{1mm}
!-------------------- Heat transport equations --------------------------
\\ \hspace*{1mm}
PET=-PEI;~~~~PE=POH+PEBM-PET*TI-PRADX;
\\ \hspace*{1mm}
PIT=-PEI;~~~~PI=PIBM-PIT*TE;
\\ \hspace*{1mm}
TE:EQ;~~~~~~~TE=TEX;~~~~~~~~TEB=TEXB;
\\ \hspace*{1mm}
TI:EQ;~~~~~~~TI=TIX;~~~~~~~~TIB=TIXB;
\\ \hspace*{1mm}
!-------------------- Poloidal field equation ---------------------------
\\ \hspace*{1mm}
HC=HCKIM;~~~~DC=DCKIM;~~~~~~XC=XCKIM;
\\ \hspace*{1mm}
CU:EQ;~~~~~~~CU=FPR;~~~~~~~~CC=CNHR;~~~~~CD=CUBM;
\\ \hspace*{1mm}
!-------------------- Additional time step control ----------------------
\\ \hspace*{1mm}
CAR27=RTOR/LTI-CAR19;~~~~~~~TSCTRL(CAR27,CAR23,CAR24,CF4):;
\\ \hspace*{1mm}
!=============== ITG transport model (end of the 1st part) ==============
\medskip\\ }
The beginning of the model is similar to already discussed.
It prescribes different density components and additional NBI heating.
Then the radial electric field is calculated which will be later used 
for computing shear of the plasma rotation velocity \twd{ROTSH}.
Two velocity components are calculated by the model. 
These are the diamagnetic velocity \twd{VDIA} and the poloidal velocity
\twd{VPOL} taken from the neoclassical theory \cite{SWW}.
The toroidal velocity \twd{VTOR} is assumed to be measured in the 
experiment and stored in the data file as \twd{VTORX}.
Note, however, that if the data for \twd{VTORX} are absent the code 
will work substituting \twd{VTOR} with zero.
The rotational shear \twd{ROTSH} is believed to reduce the effective
increment of the ITG instability and, consequently, the anomalous
transport. 
This effect is taken into account inside the calling subroutine
\twd{IFSPPL}. 

Because of very ``stiff'' character of the IFS/PPPL model a special
care should be taken to obtain a stable solution.
This is implemented in two lines after the subroutine call.
The subroutine returns electron and ion heat conductivities in two
Astra vectors \twd{CAR23} and \twd{CAR24}, respectively.
These two vectors are then successively processed with two subroutines 
\twd{SMOOTH} and \twd{FTAV} performing space and time averaging,
respectively, and yielding two smoothed vectors \twd{CAR17} and 
\twd{CAR18}.
As will be seen from the second part of the model, the radial output 
allows to view both original (i.e. calculated by \twd{IFSPPL}) and
smoothed quantities in the same plot and thus to control the quality of 
the procedure applied.
Adjusting the two control parameters \twd{CV1} and \twd{CV2} the user 
can regulate each type of averaging and achieve nearly the complete 
coincidence between the original and smoothed heat conductivities 
in the steady state.

On the other hand, a caution should be made against using this scheme
for fast processes such as the heat wave propagation.
If any of the characteristic length (for the space smoothing) or time
(for the time averaging) is larger than the wavelength or the period 
of the heat wave under consideration then an essential physic 
information can be lost in course of averaging and the applicability 
of results should be additionally inspected.

Another possibility of the accuracy control is provided by the Astra 
subroutine \twd{TSCTRL} \linebreak (time step control).
This subroutine checks if a relative variation in any of three 
quantities
\\
\hspace*{2em}{\tt CAR27~} being
a difference between the instability threshold and the actual gradient,
\\
\hspace*{2em}{\tt CAR23} $=\chi_e^{ITG}$ being the
electron heat conductivity as output from \twd{IFSPPL},\\
\hspace*{2em}{\tt CAR24} $=\chi_i^{ITG}$ being the
ion heat conductivity as output from \twd{IFSPPL},\\
exceeds a confidence level given by the parameter \twd{CF4}.
If yes, then the time step is automatically reduced.
Finally, the user can reduce the time step manually until a very low
value when no numerical instability takes place and check which of
the described procedures is most suitable in every particular 
case\footnote
{However, we have to warn the user that this check is possible 
only when the code is installed in double precision version. 
Otherwise, the time step cannot be reduced below a certain value
which is limited due to truncation errors in $U_{pl}=\prt\psi/\prt t.$}.
To conclude this discussion we note that this numerical instability
is not very dangerous and the simulation can safely continue also 
when a developed instability is clearly seen as oscillations in 
$\chi_{e,i}^{ITG}.$ 
Even then the electron and ion temperatures show quite quiet behavior
being hardly affected by a noise in the heat conductivity.

Other terms of heat conductivity equations and the equation for 
the poloidal field do not have any special features, therefore,
we consider now the rest of this model.\\ 
\noindent
{\tt\hspace*{1mm}
!=============== ITG instability based transport model (cont.) ==========
\\ \hspace*{1mm}
!-------------------- Auxiliary quantities ------------------------------
\\ \hspace*{1mm}
CV5=0.8E-3*VINT(PBLONB);~~~~~~~~!~Parallel and perpendicular 
\\ \hspace*{1mm}
CV6=1.6E-3*VINT(PBPERB);~~~~~~~~!~energy contents in the fast ions 
\\ \hspace*{1mm}
CAR28=BTOR*FRS*MU*VTOR/RTOR;
\\ \hspace*{1mm}
CAR29=HEXP;~~~~~~~~~SMEARR(CV3,CAR29,CAR30):;
\\ \hspace*{1mm}
CAR31=XEXP;~~~~~~~~~SMEARR(CV3,CAR31,CAR32):;
\\ \hspace*{1mm}
!====================== Profile output ==================================
\\ \hspace*{1mm}
Tex\bs TEX\bs\bs TEX\bs -5;~~~RLTi\bs RTOR/LTI\bs -6;
			~~~Hex\bs HEXP\bs 5;~~~~~Hix\bs XEXP\bs 5;
\\ \hspace*{1mm}
Tix\bs TIX\bs\bs TIX\bs -3;~~~RLTi\bs RTOR/LTI\bs -2;
			~~~KDes\bs CAR17\bs 5;~~~KDis\bs CAR18\bs 5;
\\ \hspace*{1mm}
Te\bs TE\bs -5;~~~~~~~~~~RLTC\bs CAR20\bs -6;
			~~~~~~ke\bs HE\bs 5;~~~~~~~~ki\bs XI\bs 5;
\\ \hspace*{1mm}
Ti\bs TI\bs -3;~~~~~~~~~~RLTD\bs CAR19\bs -2;
			~~~~~~KDe\bs CAR23\bs 5;~~~~KDi\bs CAR24\bs 5;
\\ \hspace*{1mm}
~~~~~~~~~~~~~~~~~Contributions to Er
\\ \hspace*{1mm}
ErVd\bs BTOR*VDIA\bs -1;~gamm\bs CAR22\bs -8;
			~~~~~~Er\bs ER;~~~~~~~~~~vpol\bs VPOL;
\\ \hspace*{1mm}
Er\bs ER\bs -1;~~~~~~~~~~Pecr\bs PEECR;
			~~~~~~~~~Pe\bs PETOT;~~~~~~~PNBe\bs PEBM;
\\ \hspace*{1mm}
ErVp\bs-VPOL*BTOR\bs -1;wExB\bs ROTSH\bs -8;
		~~~~~~shat\bs SHEAR;~~~~~vtor\bs VTOR\bs\bs VTORX;
\\ \hspace*{1mm}
ErVt\bs CAR28\bs -1;~~~~~betj\bs BETAJ;
			~~~~~~~~~Pi\bs PITOT;~~~~~~~PNBi\bs PIBM;
\\ \hspace*{1mm}
~~~~~~~~~~~~Current density
\\ \hspace*{1mm}
j\bs CU\bs -4;~~~~~~~~jIC\bs CUICR;~~~~~~~~nuis\bs NUIS;
						~~~Vtor\bs UPL;
\\ \hspace*{1mm}
jNB\bs CUBM\bs -4;~~~~Zeff\bs ZEF\bs\bs ZEFX;~~~tpf\bs TPF;
						~~~~~mu\bs MU\bs 1;
\\ \hspace*{1mm}
joh\bs CUOHM\bs -4;~~~jEC\bs CUECR;~~~~~~~~nues\bs NUES;
						~~~V||\bs ULON;
\\ \hspace*{1mm}
jBS\bs CUBS\bs -4;~~~~sigm\bs CC;~~~~~~~~~~betj\bs BETAJ;
						~~q\bs 1./MU\bs 5;
\\ \hspace*{1mm}
~~~~~~~~~~~~Heat conductivity
\\ \hspace*{1mm}
ncCH\bs HNCHI\bs 10;~~~GSba\bs HNGSB\bs 10;~~~grVt\bs grad(VTORX);
						~~~Hex\bs HEXP\bs 10;
\\ \hspace*{1mm}
ncPS\bs HNPSI\bs 10;~~~KDes\bs CAR17\bs 10;~~~KDis\bs CAR18\bs 10;
						~~~~~~prad\bs PRADX;
\\ \hspace*{1mm}
ncGS\bs HNGSI\bs 10;~~~GSpl\bs HNGSP\bs 10;~~~HeGS\bs HNGSE;
					~~~~~~~~~Hix\bs CAR32\bs 10;
\\ \hspace*{1mm}
ncfw\bs FOWC*HNCHI\bs 10;~KDe\bs CAR23\bs 10;~KDi\bs CAR24\bs 10;
						~~~~~~~PN\bs PENEU-PINEU;
\\ \hspace*{1mm}
~~~~~~~~~~~~Other quantities
\\ \hspace*{1mm}
ne\bs NE\bs -11;~~~~~~~dLT1\bs CAR27;~~~~~~RLT\bs RTOR/LTI\bs -2;
						~~~nB\bs NIBM\bs -11
\\ \hspace*{1mm}
pNBl\bs PBLON;~~~~~~RLNI\bs RTOR/LNI;~~~RLT\bs RTOR/LTI\bs -6;
						~~~nB1\bs NNBM1
\\ \hspace*{1mm}
ni\bs NI\bs -11;~~~~~~~shat\bs SHEAR\bs 2.;~~~RLTD\bs CAR19\bs -2;
					~~~~~nD\bs NDEUT\bs -11
\\ \hspace*{1mm}
pNBp\bs PBPER;~~~~~~RLNE\bs RTOR/LNE;~~~RLTC\bs CAR20\bs -6;
					~~~~~Zeff\bs ZEF\bs\bs ZEFX
\\ \hspace*{1mm}
!====================== Time output =====================================
\\ \hspace*{1mm}
Tix0\us TIXC\us -2;~~~~Ipl\us IPL;~~~~~~~~~Tex0\us TEXC\us -1;
					~~<Te>\us TEAVB;
\\ \hspace*{1mm}
Ti0\us TIC\us -2;~~~~~~PNBI\us QNBI;~~~~~~~Te0\us TEC\us -1;
					~~~~<Ti>\us TIAVB;
\smallskip\\ \hspace*{1mm}
ne0\us NEC;~~~~~~~~~Wthm\us WTOTB;~~~~~~<ne>\us NEAVB;
					Wtot\us WTOTB+CV5+CV6;
\\ \hspace*{1mm}
nlc\us NECHC;~~~~~~~betp\us BETAJB;~~~~~Zef0\us ZEFC;
					~Wequ\us WTOTB+0.75*CV6+1.5*CV5
\smallskip\\ \hspace*{1mm}
TiB\us TIB;~~~~~~~~~Ipl\us IPL;~~~~~~~~~iter\us TITER;
					~~~~INB\us IBMB;
\\ \hspace*{1mm}
TeB\us TEB;~~~~~~~~~Ulc\us UPLB;~~~~~~~~tauE\us TAUEB;
					~~~~IBS\us IBSB;
\smallskip\\ \hspace*{1mm}
Ptot\us QTOTB;~~~~~~POH\us QOHB;~~~~~~~~Pe\us QETOTB;
					~~~~~PNBA\us QBTOTB;
\\ \hspace*{1mm}
Pi\us QITOTB;~~~~~~~PNB\us QNBI;~~~~~~~~shin\us QNBI-QBTOTB;
					~Prad\us QRADB;
\medskip\\ }
{\tt\hspace*{1mm}
!======================== ITG model end =================================
}\vskip 1ex
After calculation of few auxiliary quantities which are used for 
output only the model defines 80 vector and 32 scalar output signals.
The signals are ordered in a way that groups of associated quantities
are placed as far as possible in the same box and on the same screen.
Some of signals are repeated that allows to present them in appropriate
context.

Measurable quantities (as $T_e$ or $T_i$) are plotted in the same box
with the calculated ones.
If corresponding experimental data are not available it is not 
considered as an error in a command like \twd{Tex\bs TEX\bs\bs TEX}.
However, Astra vectors \twd{XEXP} and \twd{HEXP} which show 
``effective'' heat conductivities include gradients of measured
quantities in a denominator that can cause division by zero.

The second group of 16 radial curves includes information concerning
the plasma rotation.
The second box here plots the ITG instability increment \twd{CAR22}
together with the rotational shear \twd{ROTSH} 
(see page \pageref{ROTSH}).
Other radial output screens present current related quantities,
different contributions to the heat conductivity, density 
distributions and some gradients of plasma parameters which are
essential for the ITG stability.

Different local and global quantities are shown as time dependent
signals.
One can see here central and edge temperatures, energy contents
calculated in three different ways, line and volume averaged
densities, energy confinement time, partial 
components of the total current and, finally, volume integrated
power sources and sinks.

\subsection{Data setting}

\subsubsection{Hierarchy of data initialization}

A chain of initialization procedures is invoked each time when the code
is started.
First of all, all variables with very few exceptions have pre-defined
(default) values. 
The default value of a variable is used only if the variable is not
defined explicitly in any other way. 
The most straightforward and natural way of data assignment is using a
start (or data) file. 
A name of this file is given by the user as a second parameter in the 
Astra command line (see Section~5.5.2).
Another possibility is assignment instruction in a model as 
described in the previous section.
One more option is an interactive change during the run time.
A sequence of assignments for different types of variables is shown in
the following table. 
\medskip\\
\indent
{\bf Table 5.1. Order of variable setting}
\nopagebreak
\\[10pt]
\nopagebreak
\noindent
\begin{tabular}{| l | c | c | c | c | c | c |}			\hline
\bf Variable type	&\bf Default	&\bf{\tt .log }file
					&\bf Data file
						&\bf Model
							&\bf Run time \\ 
								\hline
Device/plasma parameters	& $+$	& $+$	& $+$	& $+$	& $+$ \\
Adjustable (dummy) variables	& $+$	& $+$	& $-$	& $+$	& $\pm$\\
Grid/time/accuracy control	& $+$	& $+$	& $-$	& $+$	& $+$ \\
Color control			& $+$	& $+$	& $-$	& $-$	& $+$ \\ 
								\hline
\end{tabular}
\\[10pt]
Every assignment type overrides all others which have 
more left position in this table.
For instance, if a quantity is redefined in the run time then every
other definitions are disabled.
The symbol ``$-$'' means that this type of variable cannot be assigned 
in this way.
The marking $\pm$ will be discussed in Section~5.3.4.

There are few exceptions from the table above.
Namely the following quantities are used for the radial grid allocation
and, therefore, they have to be defined at earlier phase of 
\linebreak\vspace*{-20pt}
\begin{tabbing}
\indent	\tt AWALL\quad  \=\quad (none)\quad \= $\rho$--grid size is 
     $\rho_B=\sqrt{\mbox{\tt AWALL*AWALL*{\rm max}(ELONM,ELONG)}},$\kill
\indent	\tt AB	  	\>(none)
			\> maximum minor radius in a horizontal plane,\\
\indent	\tt RTOR  	\>(none)
			\> major radius at the same plane,	\\
\indent	\tt ELONM 	\>(1)	
			\> maximum elongation,			\\
\indent	\tt TRICH 	\>(0)	
			\> chamber triangularity (used for drawing
			   only),\\
\indent	\tt AWALL 	\>\tt (AB)
			\> $\rho$--grid size is 
     $\!\!\rho_W=1.1\tw{*AWALL*}\sqrt{\mbox{\rm max\tt(ELONM,ELONG)}},$\\
\indent	\tt NB1	  	\>(41)
		\> $\rho$--grid step is   $h=\rho_W/(\mbox{\tt NB1}-0.5)$
\end{tabbing}
\noindent
the code execution and cannot be redefined later. 
It means that these six quantities can be defined either by a
\tw{`.log'} file (see Section 5.3.2) or by a data file and 
they cannot be set in a model or changed during the run.
This limitation seems to be not very restrictive while it arises 
quite naturally from the requirement
that these basic variables describing geometry of a particular
device cannot depend on time. 
When applicable, the quantities acquire default values which are 
shown in brackets.
However, it is recommended that the first 
four of these quantities are defined explicitly. 
In opposite, the default definition for {\tt AWALL} is in most cases
well sufficient, therefore, it should be defined by the user only if 
Astra provides a diagnostic that the default value is not appropriate.
When setting the grid size {\tt NB1}, the user should take into account
that the grid size (variable {\tt NA1} which will be automatically
selected by the code) actually used in a simulation is usually a
fraction ($\geq$80\%) of \twd{NB1}. 
The rest of the radial grid is reserved for a description of the
scrape-off layer and for a possible increase of the Lagrangian variable
$\rho_B={\tt RHO(NA1)}<\rho_W.$ 

We continue a list of variables which should be defined by the user.
In addition to already mentioned \twd{AB, RTOR, ELONM, TRICH,} those are
(default values are shown in parentheses)
\\
\centerline
{\tt ABC(AB), SHIFT(0), ELONG(ELONM), TRIAN(0), BTOR{\rm(none)}, ZMJ(1),
AMJ(2).}
\\
Finally, at least one of the three quantities
\\
\centerline{\tt
IPL{\rm(none)}, UEXT{\rm (none)}, LEXT\rm(none)}
\\
also must be defined in order to avoid a possible crash during the code
start-up. 
However, all these quantities can be defined by every method listed in
the first line of Table~5.1.

\paragraph{Default array setting.}

Unlike simple variables arrays can be defined either via a data file or
in a model.
Following is the list of vector variables (arrays) which can aquire
default definitions 
\begin{tabbing}
$\tw{NE}(\rho)=\tw{NEX}({\rho}),$\quad	\=
$\tw{F1}(\rho)=\tw{F1X}(\rho),$\quad	\=
$\tw{NI}(\rho)=\tw{NE}(\rho)/\tw{ZMJ},$\quad\=
$\tw{CU}(\rho)\propto\tw{CC}(\rho)$\kill
$\tw{NE}(\rho)=\tw{NEX}(\rho),$	\>
$\tw{F1}(\rho)=\tw{F1X}(\rho),$	\>
$\tw{AMAIN}(\rho)=\tw{AMJ},$	\>
$\tw{CC}(\rho)=\tw{CCSP}(\rho),$\\
$\tw{TE}(\rho)=\tw{TEX}(\rho),$	\>
$\tw{F2}(\rho)=\tw{F2X}(\rho),$	\>
$\tw{ZMAIN}(\rho)=\tw{ZMJ},$	\>
$\tw{CU}(\rho)\propto\tw{CC}(\rho),$\\
$\tw{TI}(\rho)=\tw{TIX}(\rho),$	\>
$\tw{F3}(\rho)=\tw{F3X}(\rho),$	\>
$\tw{NI}(\rho)=\tw{NE}(\rho)/\tw{ZMJ},$	\>
$\tw{ZEF}(\rho)=\max(1.,\tw{ZEFX}(\rho)).$
\end{tabbing}
These definitions are effective only if the left hand sides are not
defined explicitly in a model. 
For instance, if \twd{NE} appears on the left hand side in an assignment
instruction in a model then the default definition \twd{NE=NEX} is not
involved. 
On the other hand, if any of the quantities {\tt NE, TE, TI} is not
assigned in a model then it must be defined in the data file.  
Otherwise, a run time error could be encountered.
This requirement is not applied to the quantities \twd{F1, F2, F3}.
Because their usage is optional they need not be defined unless they are
directly used in a model.

\subsubsection{Reading from a ``{\tt.log}'' file}

Usually a large amount of tuning parameters participate in a transport
modelling. 
Often these parameters are not known in advance, therefore, an
adjustment in course of simulation is a most reasonable way for finding
appropriate values. 
The Astra code provides a possibility to store all internal
variables adjusted during the current run.
This happens when user selects the option ``Save const'' (or the hot key
\twd{`I'}) from the run time menu (Section~5.5.3).
The current values of adjustable variables are saved in a file
which is read each time when the code starts up so that all variable 
values tailored and saved in a previous simulation will be used in the
subsequent one. 
A name of the file is composed by adding the extension ``{\tt .log}'' 
to a name of the running model. 
This convention implies that only one {\tt .log}~~file, and,
consequently, one set of stored values corresponds to each model file. 
These files can also be edited manually.

\subsubsection{Input from a data file}

Unlike the ``{\tt .log}'' file which belongs to a corresponding model
there can be an arbitrary number of input files which describe an
experimental setup and are not related to any particular model.
Reading input data from a file is a conventional way of data definition. 
The name of this file is specified as a first parameter
in a command line starting execution of the Astra code.
The format of this file is described in Section~5.4. 
Here we discuss only the main principles of data input by making
use of the data file.

First of all, note that starting from this level of data initialization
(see Table~5.1) every scalar or vector variable can be defined not only
as a time constant but also as a time dependent function.
On one hand, this adds a flexibility to the data setting. 
On the other hand, it results in a restriction that only a limited
number of specially treated variables can be read from a data file. 

\paragraph{Scalar variables readable from a data file.}

We start with a note that in future versions of the Astra code the set
simple variables readable from a data file may be changed but the 
updated list (except for the group \twd{ZRD1, ZRD2,}\dots and the three
variables \twd{NB1, TSTART, TEND\rm)}
can always be viewed during the code run by pressing the key `V' or 
by clicking mouse in the field ``Variables''.
In the version 5.3 of the code this set comprises
\medskip\\
\indent
{\bf Table 5.2. List of Astra scalars definable by a data file}
\nopagebreak
\\[10pt]
\nopagebreak
\noindent
{\tt
\begin{tabular}{|p{15mm}p{15mm}p{15mm}p{15mm}p{15mm}p{15mm}p{15mm}p{15mm}|}
\hline
AB	&ABC	&AIM1	&AIM2	&AIM3	&AMJ	&AWALL	&BTOR	\\
ELONG	&ELONM	&ENCL	&ENWM	&FECR	&FFW	&FICR	&FLH	\\
GN2E	&GN2I	&IPL	&LEXT	&NNCL	&NNWM	&QECR	&QFW	\\
QICR	&QLH	&QNBI	&RTOR	&SHIFT	&TRIAN	&TRICH	&UEXT	\\
UPDWN	&WNE	&WTE	&WTI	&ZMJ	&NB1	&TSTART	&TEND	\\
ZRD1	&ZRD2	&ZRD3	&ZRD4	&ZRD5	&ZRD6	&ZRD7	&ZRD8	\\
ZRD9	&ZRD10	&ZRD11	&ZRD12	&ZRD13	&ZRD14	&ZRD15	&ZRD16	\\
ZRD17	&ZRD18	&ZRD19	&ZRD20	&ZRD21	&ZRD22	&ZRD23	&ZRD24	\\
ZRD25	&ZRD26	&ZRD27	&ZRD28	&ZRD29	&ZRD30	&ZRD31	&ZRD32	\\
ZRD33	&ZRD34	&ZRD35	&ZRD36	&ZRD37	&ZRD38	&ZRD39	&ZRD40	\\
ZRD41	&ZRD42	&ZRD43	&ZRD44	&ZRD45	&ZRD46	&ZRD47	&ZRD48	\\
\hline
\end{tabular}
}\\[10pt]
Each variable from this list is mirrored with another variable with 
a similar name but with trailing ``{\tt X}''(which stands for
eXperimental), X-scalars.
\medskip\\
\indent
{\bf Table 5.3. List of X-scalars readable from a data file}
\nopagebreak
\\[10pt]
\nopagebreak
\noindent
{\tt
\begin{tabular}{|p{15mm}p{15mm}p{15mm}p{15mm}p{15mm}p{15mm}p{15mm}p{15mm}|}
\hline
ABX	&ABCX	&AIM1X	&AIM2X	&AIM3X	&AMJX	&AWALLX	&BTORX	\\
ELONGX	&ELONMX	&ENCLX	&ENWMX	&FECRX	&FFWX	&FICRX	&FLHX	\\
GN2EX	&GN2IX	&IPLX	&LEXTX	&NNCLX	&NNWMX	&QECRX	&QFWX	\\
QICRX	&QLHX	&QNBIX	&RTORX	&SHIFTX	&TRIANX	&TRICHX	&UEXTX	\\
UPDWNX	&WNEX	&WTEX	&WTIX	&ZMJX	&	&	&	\\
ZRD1X	&ZRD2X	&ZRD3X	&ZRD4X	&ZRD5X	&ZRD6X	&ZRD7X	&ZRD8X	\\
ZRD9X	&ZRD10X	&ZRD11X	&ZRD12X	&ZRD13X	&ZRD14X	&ZRD15X	&ZRD16X	\\
ZRD17X	&ZRD18X	&ZRD19X	&ZRD20X	&ZRD21X	&ZRD22X	&ZRD23X	&ZRD24X	\\
ZRD25X	&ZRD26X	&ZRD27X	&ZRD28X	&ZRD29X	&ZRD30X	&ZRD31X	&ZRD32X	\\
ZRD33X	&ZRD34X	&ZRD35X	&ZRD36X	&ZRD37X	&ZRD38X	&ZRD39X	&ZRD40X	\\
ZRD41X	&ZRD42X	&ZRD43X	&ZRD44X	&ZRD45X	&ZRD46X	&ZRD47X	&ZRD48X	\\
\hline
\end{tabular}
}\\[10pt]
In particular, it means that there are couples of independent variables,
\twd{AB} and \twd{ABX}, \twd{ABC} and \twd{ABCX} and so on, which can be
defined and used quite independently with one significant exception.
Namely, if any of quantities from the second list is not defined
explicitly in a model then, by default, it takes value of its
experimental prototype. 
For a practical work it means that if the quantity \twd{IPLX}
is defined by a data file and the quantity \twd{IPL} is not
mentioned in a model then there is no difference between both.
In this case, the resulting code is fully equivalent to the code
obtained with the explicit instruction \twd{IPL=IPLX}.
However, if the quantity \twd{IPL} is defined in some independent way
(for instance, by the boundary condition \twd{UEXT=0}) then 
\twd{IPLX} and \twd{IPL} are fully independent. 

At first sight, this convention could seem too complicated but
it works so that a novice user of the code need not know anything
about two different sets of data. 
This allows to omit in a model multiple instructions of type 
\twd{ABC=ABCX}.
However, an advanced user can take advantage of that keeping trace
on \twd{IPLX} as defined by experimental measurements and comparing
it with \twd{IPL} calculated in the code.
Note also, that if more complicated condition for \twd{IPL} is used,
for instance, if \twd{IPL} is calculated inside user's subroutine 
then the Astra code builder cannot recognize this hidden definition
and will try to set \twd{IPL=IPLX} in conflict with user's intention.
A simple way to overcome this contradiction is using a dummy
instruction \twd{IPL=IPL} in a model.

Three Astra control variables \twd{NB1, TSTART} and \twd{TEND} are used
for influencing a simulation flow. 
They have no X-counterpart, nevertheless, it is often useful to define
these quantities by a data file. 
Therefore they are added to the list of variable in Table~5.2.

\paragraph{Vector variables readable from a data file.}
Only a limited number of basic vector variables (\tw{X}-vectors) can be
defined by an input file. 
\medskip\\
\indent
{\bf Table 5.4. List of X-vectors readable from a data file}
\nopagebreak
\medskip\\
\nopagebreak
\noindent
{\tt\begin{tabular}
{|p{15mm}p{15mm}p{15mm}p{15mm}p{15mm}p{15mm}p{15mm}p{15mm}|}
\hline
TEX	&TIX	&NEX	&NIX	&CUX	&ZEFX	&VPOLX	&VTORX 	\\
MUX	&MVX	&GNX	&SNX	&PEX	&PIX	&PRADX	&IPOLX	\\
SHX	&ELX	&TRX	&VRX	&G11X	&G22X	&G33X	&DRODAX	\\
F1X	&F2X	&F3X	&CAR1X	&CAR2X	&CAR3X	&CAR4X	&CAR5X	\\
CAR6X	&CAR7X	&CAR8X	&CAR9X	&CAR10X	&CAR11X	&CAR12X	&CAR13X	\\
CAR14X	&CAR15X	&CAR16X	&	&	&	&	&	\\
\hline
\end{tabular}
}\bigskip\\
The list includes 
\begin{description}
\vspace*{-1ex}\item[--] 
the main plasma characteristics which can be measured in an experiment 
(e.g. electron temperature \twd{TEX}, density \twd{NEX,} effective 
charge \twd{ZEFX} and so on), 
\vspace*{-1ex}\item[--]
calculated by external codes (e.g. quantities supplied by 
equilibrium, heating and current drive codes), 
\vspace*{-1ex}\item[--] 
a set of 19 dummy arrays  \twd{(F1X, F2X, F3X, CAR1X, ..., CAR16X)}
which can acquire any meaning depending on user's will. 
\end{description}\vspace*{-1ex}
Any array of Table~5.4 will be defined if its name appears in a data
file according to the rules discussed in Section~5.4.3 and this data
file is used for the current run.

As already discussed in Section~5.3.1, a default definition of the type 
$\tw{NE}(\rho)=\tw{NEX}({\rho})$ is applied to the following seven
\tw{X}-arrays \twd{NEX, TEX, TIX, F1X, F2X, F3X, ZEFX} only. 
The reason is that undefined \twd{NE, TE, TI, ZEF} can cause a
simulation crash and \twd{F1, F2, F3} are added in order to keep
similarity between the main and auxiliary transport equations, \req{main}
and \req{aux}, respectively. 

Usually, the default assignments are effective in absence of
corresponding explicit definition only.
However, there is one exception.
If the equilibrium control parameter \twd{NEQUIL} (default value 0) is
set as \\
\hspace*{2cm}\tw{NEQUIL=-1}\\
then it is assumed that all the geometry and configuration related
quantities are defined by means of a data file. 
In this case, nine \tw{X}-arrays \twd{SNX, ELX, TRX, VRX, DRODAX, IPOLX, 
G11X, G22X, G33X} must be pre-calculated and stored in the data file.  
These nine quantities will then be used to calculate the corresponding
set of Astra vector variables as
\begin{tabbing}
$\tw{VR}(\rho)=4\pi^2\rho\,\tw{*RTOR*VRX}(\rho)$	\=
$\tw{ELON}(\rho)=\tw{ELONG*ELX}(\rho)$			\=
$\tw{TRIA}(\rho)=\tw{TRIAN*TRX*}(\rho/\rho_a)^2,$	\kill
$\tw{SHIF}(\rho)=\tw{SHIFT*SHX}(\rho),$			\>
$\tw{ELON}(\rho)=\tw{ELONG*ELX}(\rho),$			\>
$\tw{TRIA}(\rho)=\tw{TRIAN*TRX*}(\rho/\rho_a)^2,$	\\
$\tw{VR}(\rho)=4\pi^2\rho\,\tw{*RTOR*VRX}(\rho),$	\>
$\tw{DRODA}(\rho)=\tw{DRODAX}(\rho),$			\>
$\tw{IPOL}(\rho)=\tw{IPOLX}(\rho),$			\\
$\tw{G11}(\rho)=\tw{VRX}(\rho)\tw{*G11X}(\rho),$	\>
$\tw{G22}(\rho)=\rho\,\tw{*G22X}(\rho),$	\>
$\tw{G33}(\rho)=\tw{G33X}(\rho).$
\end{tabbing}
As one can see, these \tw{X}-arrays are selected in such a way
that for concentric non-shifted circle magnetic surfaces they all are
unitary. 
These assignments are always operative when $\tw{NEQUIL}=-1$ and cannot
be overwritten by any assignments in a model.
On the contrary, when $\tt NEQUIL\neq -1$ these nine \tw{X}-arrays have
nothing to do with the corresponding left hand sides and can be used for
any other goals.

Similarly, all other \tw{X}-arrays have no pre-defined meaning and can
be used for input of any radial profile. 
For instance, if the vacuum rotational transform \twd{MV} should be
taken from a data file then the assignment \twd{MV=MVX} has to be
made explicitly in a model. 
Without this assignment \twd{MVX} can be anything.
The \tw{X}-arrays can be used in expressions according to usual rules
\\
\hspace*{2cm}\tw{PE=PEX-PRADX+POH-PEICL}
\\
or appear on the left hand side of an assignment command,
e.g., \twd{CAR1X=FPR**1.5}.
However, the latter construction is not recommended because if the same
array \twd{CAR1X} is occasionally defined in a data file then a result is
unpredictable. 

\subsubsection{Run time variable setting}

At the run time most of variables are accessible for a direct
control by the user. 
There are two ways of doing this.
First of all, it can be done interactively by pressing keys 
`A', `C', `D' and `V' or by mouse selecting proper fields from the 
run time menu. 
Another way is defining some quantities in user's subroutine.
This method can be used when a quantity has to be calculated on a
basis of complicated formula or equation.

As follows from Table 5.1, definition of any quantity during the code
execution has the highest priority and overrides any other definitions.
In particular, once any plasma or device parameter is interactively
changed by the user its previous assignments in a data file or in a
model are discarded. 

Unfortunately, it is more difficult to fulfill the same property for a 
C-parameter.
One should beware that a result of interactive setting C-parameter can
be guaranteed only if it is not assigned in a model. 
If, nevertheless, a C-parameter is defined in the model by a
time-independent instruction like \tw{CV1=2} then a run-time
modification will work properly, however, it will not in case of more
complicated time-dependent assignment as \tw{CV1=FJUMP(0.1)}. 

A similar difficulty arises when a variable or C-parameter is defined in
a user subroutine. 
If it is simultaneously defined in a model or changed interactively then
a result is unpredictable.
Therefore, it is recommended to avoid any conflicting definitions
of those types.

\subsection{Data file format}

During installation of the Astra system the user is supplied with 
sample data files. 
New data files can be obtained by editing the 
original files retaining their structure. 
Therefore, this section can be omitted during the first reading. 
For more elaborated data setting we describe syntax of a 
data file in this section.

\subsubsection{General requirements}

In what follows we will distinguish between input of simple variables 
and vectors. 
The former means that a scalar function of time will be defined by the
input. 
The latter defines a time dependent vector function of a magnetic
surface while the ``radial'' coordinate can be selected by the user from a
wide set of different options. 
It is also possible to define this function using two spatial
coordinates $(r,z)$ in the poloidal plane.  
This 2D function is then mapped to a 1D function of magnetic surfaces 
using the time dependent magnetic configuration of the particular Astra
run. 

The first two lines in a data file are reserved for a commentary. 
They are not used by the code itself but the first 16 positions of 
the first line appear in the output graphic window during the Astra run.
For instance, they can be used for displaying a device name and a shot 
number under consideration.
A length of meaningful string in a data file cannot exceed 132
characters. 
There is no limit for a length of commentary string.

During reading a data file each string is checked whether its first 6
characters match one of keywords. 
The keywords are all the names listed in Tables 5.3 or 5.4
and, additionally, the following seven control words: 
\\ 
{\tt\indent
NAMEXP, NTIMES(1), FACTOR(1), FILTER(0.001), GRIDTYPE, POINTS, END.
}\\
A variable corresponds to each of the first 6 control words.
An input value of this variable should follow after the control word.
Once defined these values are valid until they are redefined by 
another control string.
The control variables \tw{FACTOR} (being the unit conversion factor),
\tw{NTIMES} (number of time slices) and \tw{FILTER} (defining a
transformation rule), if not defined explicitly are set to a default
value shown in parentheses. 
When the control word \tw{END} is found then reading the data file 
is ceased, otherwise, it is continued until the file end.
A meaning of other variables will be discussed later.
All names and control words can appear in low or upper case.

If a string is started with one of the control words it
is treated as {\bf a control string}.
When a control string is encountered it is parsed and an interpretation
regime for a group of several subsequent lines is set.
When the group is exhausted then the following 
string is again checked against the matching condition. 
If first 6 characters of the string do not match any of the keywords 
then the string is skipped and the next string is processed similarly. 
For instance, any non-alphabetic symbol in the first position of 
a parsed string results in ignoring the string by the code.
This property can be used for writing commentaries in the input file.

For a sake of brevity the character \twd{`X'} on the end of 
every name in the data file can be omitted. 
Then the \twd{`X'} is added by the code so that presence or absence
of \twd{`X'}in a name does not change the result. 
For instance, there is no matter whether a name from Table 5.3 or 
Table 5.2 is used. 
But in any case, a data file defines a variable from Table 5.3 (or 5.4)
only. 
Whether the corresponding variable without the trailing \twd{`X'} is
defined or not depends on the contents of a model file and
interpretation rules discussed in the previous section.

When reading input data is started the code is in a scalar input mode. 
It means that at first all scalar variables should be defined. 
After one of the control words
{\tt~POINTS, GRIDTYPE, FILTER~} is encountered the code is switched
to the vector input mode.
Now the vector variables (Table 5.4) have to be defined.
If a simple variable appears on the input now it will be ignored.

There are several different ways of data input in the Astra system. 
It can be either a direct reading data from a start file or redirecting
input to so called ``U-file'' where the data are represented in a U-file
standard developed in PPPL for storage of experimental data. 
The U-file standard is used also for the ITER data base. 
This standard is not described in this paper but it is quite
transparent and can be easily understood from few examples attached
to the Astra code.
%Recently, a facility to read so call ``Ex-file'' was added. 
%Ex-file is a standard used in the JETTO code \cite{JETTO}.
%Reading Ex-files requires libraries available at JET only, therefore, it
%is not discussed here.

Different input formats can be combined in one input file.
Moreover, mixing different formats is permitted for different time
slices of the same input quantity.
Definition of a time dependent scalar or vector can consist of several
subsequent groups for different time slices. 
However, one restriction should be mentioned here. 
These groups can use different grids but the
subsequent times should be given in the increasing order and
all groups related to the same variable should be contiguous.
In other words, when a time evolution for one variable is defined 
and the input sequence is switched to another variable 
then return to the first one is forbidden.

\subsubsection{Reading data from a U-file.}

First of all, the first 6 characters of every line are considered.
If they are recognized as a name of variable then the next word in the
string is analyzed. 
If this word is \twd{"U-file"} (case insensitive) then the string is
interpreted as a reference (pointer) to a U-file where a data set for 
the current variable is stored. 
The name of the U-file should be separated from the word \twd{U-file}
with a colon  \twd{`:'}.
Optionally, a multiplicative factor can be given after the 
second~\twd{`:'}. 

The following 5 lines show different allowed formats of a pointer to
U-file: 
\\
\dl{IPL \tab U-file\tab:\tab A00000.aaa\tab Factor:\tab 0.000001}
\dl{IPL \tab u-file\tab:A00000.aaa\tab factor:1.E-6}
\dl{IPLX\tab U-FILE\tab:A00000.aaa\tab:1.E-6}
\dl{iplx\tab U-File\tab:A00000.aaa:1.E-6}
\dl{ipl \tab u-file:A00000.aaa:1.E-6}
\\
All these lines have the same effect and mean that the variable
describing the plasma current, \twd{IPLX}, should be taken from the
U-file \twd{AWD/udb/A00000.aaa} (\twd{`AWD'} denotes the Astra working
directory) where it is stored in amperes and multiplied by 
$10^{-6}$ because the Astra code employs MA as a current unit. 
The word \twd{"Factor"} can be omitted.
The symbol \tab~ shows a delimiter, i.e. any combination of 
spaces and horizontal tabulations.

Note that for this input format
\\ -- only a U-file name is case sensitive,
\\ -- only one line is allowed for every variable,
\\ -- the word ``factor'' and the factor value can be omitted,
\\ -- if the factor value is omitted it is set to 1,
\\ -- a delimiter ``\tab'' is required only
\\ \hspace*{5mm} -- between a name of variable and the word ``U-file'',
\\ \hspace*{5mm} -- between a name of U-file and the word ``factor'' 
if the latter is present. 

The U-file name should be given with a path relative to the directory
\twd{AWD/udb/}. 
For instance, if the name is \twd{../exp/A00000.aaa } then the
U-file containing data should be in the subdirectory \twd{AWD/exp/}.
Every permissible U-file name is allowed but the file must exist and
contain data. 
For input of a scalar variable it should be a 0D or 1D U-file. 
Moreover, in the latter case it should have time as independent variable.
For input of a radial array it should be 1D or 2D U-file.

In the ITER data base every U-file consists of a number of merged
separate U-files for different signals.
In order to select an appropriate record one can use a specifier
\twd{`<'} as shown below
\\
\dl{IPL~~~U-file:JET/jet\us 40847\us 1D~<~IP~~:1.E-6}
\dl{ABC~~~U-file:JET/jet\us 40847\us 1D~<~AMIN:1.E-6}
\dl{BTOR~~U-file:JET/jet\us 40847\us 1D~<~BT~~:1.E-6}
\dl{ZRD1~~U-file:JET/jet\us 40847\us 1D~<~LI~~:1.E-6}
\\
In this example, all four Astra variables \twd{IPLX}, \twd{ABCX},
\twd{BTORX} and \twd{ZRD1X} are read from the same U-file 
\twd{AWD/udb/JET/jet\us 40847} 
where they are stored under the names \twd{IP}, \twd{AMIN}, \twd{BT} and
\twd{LI}, respectively. 

We also remind that, first of all, scalar variables have to be defined
and then the input mode can be switched to the vector input.
Thus, the group
\\
\dl{IPL \tab u-file:A00000.aaa:1.E-6}
\dl{FILTER\tab 0.002}
\dl{TEX \tab u-file:B00000.bbb:factor\tab 1.E-3}
\dl{MV~  \tab u-file:C00000.ccc}
\\
first defines the simple variable \twd{IPLX} as discussed before.
Then the control word \tw{FILTER} switches the input regime to the
vector mode.
(It could also be the words \twd{POINTS} or \twd{GRIDTYPE} but these two
words do not affect the input from a U-file.)
Then the radial array \twd{TEX} is taken from the U-file 
\twd{AWD/udb/B00000.bbb} and all array values are multiplied by the
factor $10^{-3}.$ 
Finally, the array \twd{MVX} will be defined.

\subsubsection{Fixed format input}

If the first word in a string is recognized as a name from Table 5.3 or
Table~5.4 and the second word is not \twd{"U-file"} then the rest of the
string will be parsed in {\bf positional or fixed format} mode because all symbols will be interpreted according to their
positions in the string. 
This format was used in old versions (before 3.0) of the Astra code. 
Although it is quite restrictive, for compatibility it is also supported
in more recent versions. 

\paragraph{Fixed format for a scalar input.}

In the scalar input mode all positions in a string have the meaning
specified by the table: 
\bigskip\\
\begin{tabular}{| l | l l l |}\hline
Field position	&1-6		& 9-14	& 17-22	\\ \hline
Field meaning	&\twd{NAME}	&Time	&Value	\\ \hline
\end{tabular}
\bigskip\\
\twd{NAME} is one of the names listed in Table 5.2 (or Table 5.3). 
``Time'' and ``Value'' 
are numbers which obey Fortran rules for integer or real constants and 
cannot use a place beyond the reserved fields of 6 positions. 
Non-used positions (except for those on the end of a string) within 
the reserved fields and the positions 7, 8, 15 and 16 should be filled 
with spaces. 
In the fixed format, the tabulations cannot be used as a separator 
between meaningful fields. 
For a time-independent variable the second field ``Time'' can be 
skipped (filled with spaces).
For instance, the string (spaces are shown with a symbol '{\tt\sp}')
\\ \dl{BTOR\sp\sp\sp\sp\sp\sp\sp\sp\sp\sp\sp\sp3.5}
is equivalent to the Fortran statement 
\\
\fl{BTORX=3.5}
A variable evolving in time can be given as
\\
\dl{Name\sp\sp\sp\sp Time\sp\sp\sp\sp Value}
\dl{123456\sp\sp 901234\sp\sp 789012}
\dl{IPL\sp\sp\sp\sp\sp .0\sp\sp\sp\sp\sp\sp.1}
\dl{IPL\sp\sp\sp\sp\sp 1.\sp\sp\sp\sp\sp\sp 1}
\dl{IPL\sp\sp\sp\sp\sp 2\sp\sp\sp\sp\sp\sp\sp 1.0}
\dl{IPL\sp\sp\sp\sp\sp 2.5\sp\sp\sp\sp\sp\sp 2.E-1}
\\
The first two lines will be interpreted by the Astra code as comment
strings.  
The whole group determines the internal variable \twd{IPLX} as a 
function of time
\begin{equation}\label{IPLX}
\twd{IPLX}(t)=\left\{
\begin{array}{ll}
0.1\QQUAD	&{\rm if}\quad t\le 0\\
0.1+0.9t	&{\rm if}\quad 0\le t\le 1\\
1		&{\rm if}\quad 1\le t\le 2\\
1-0.8(t-2)	&{\rm if}\quad 2\le t\le 2.5\\
0.2		&{\rm if}\quad t\geq 2.5\\
\end{array}\right.
\end{equation}
As can be easily seen the time evolution is obtained by a linear
interpolation between subsequent times. 
Outside the specified time interval the quantity is extended as
a constant.

\paragraph{Fixed format for arrays.} 
As already mentioned, the Astra code starts reading data file being
in the scalar input mode. 
In order to invoke vector mode the variable \twd{POINTS} must be
defined. 
It can be done by a line as
\medskip\\
\dl{POINTS\tab\tt 20}
\\[1ex]
This line switches the input mode to reading vector variables and
additionally declares that the vector dimensionality is 20.
Then the first 6 characters in a line will be compared with array names
of Table~5.4.
If it fits one of those name and the next word is not \twd{"U-file"} 
then the string should have the following fixed format:
\bigskip\\
\begin{tabular}{| l | c c c c c c c|}\hline
Field position	&1-6	&7-11	&12-16	&17-21	
&\dots		&\tt 7+5*j-11+5*j &\dots\\ \hline
Field meaning	&\twd{NAME}& Time& Value& Value	
&\dots		& Value	      &\dots\\ \hline
\end{tabular}
\bigskip\\
All positions which are not used must be filled with spaces. 
On the other hand, no spaces is required between the different fields. 
The total number of fields in a line is $\twd{POINTS}+2$ and 
$1\le j\le\tt POINTS.$ 

As in the case of scalar variables the first two fields in the line give
a variable name and a time value.
If the variable is time independent then the second field can be empty
(filled with spaces).
Remaining fields define radial dependence of the vector assuming that
its values are given on the equidistant grid in a radial variable.
By default this radial variable is $0\leq a\leq a_B=\tt ABC.$
This default convention can be changed by redetermining variable
\twd{XINPUT} (see Section~4.6.1). 

A time dependent vector can be defined similar to a time dependent
scalar. 
In the example below, the time evolving density \twd{NEX} is determined. 
\bigskip\\
\dl{POINTS\tab\tt 10}
\dl{123456789012345678901234567890123456789012345678901234567890}
\dl{Name\sp\sp Time\sp ValueValueValueValueValueValueValueValueValueValue}
\dl{NE\sp\sp\sp\sp .0\sp\sp\sp 2.30\sp 2.27\sp 2.23\sp 2.15\sp 2.02\sp 1.8\sp\sp 1.49\sp 1.13\sp 0.8\sp\sp .522}
\dl{NE\sp\sp\sp\sp 0.2002.3842.4302.4562.4102.3022.1281.9031.6471.3601.150}
\dl{NE\sp\sp\sp\sp .4\sp\sp\sp 2.66\sp 2.77\sp 2.85\sp 2.86\sp 2.76\sp 2.58\sp 2.33\sp 2.03\sp 1.68\sp 1.43}
\dl{TE\tab U-file:T12345.ECE:.001}
\dl{TI\sp\sp\sp\sp\sp\sp\sp\sp\sp 3.22\sp 3.14\sp 3.07\sp 2.64\sp 2.07\sp 1.57\sp 1.14\sp 0.71\sp \sp .25\sp\sp.12}
\dl{END}
The first line in this example sets value \twd{POINTS=10} for all
subsequent lines until it will be changed by another command line. 
The second and third lines label positions.
As long as they include no valid name in the first six positions they
are ignored by the code. 
As shown in the second line for ``\twd{NEX}'' no delimiter is required 
between the different fields.
Note that a U-file \twd{T12345.ECE} can have a number of radial points
different from 10 but it will not affect the current value of
\twd{POINTS.} 
Because the length of every input line is limited by 132 positions no
more than 24 array elements can be used in this format.

\subsubsection{Free format input.}

The control words {\tt~NAMEXP, NTIMES, FACTOR~} can be used for
input both scalar and vector variables.
Moreover, any of the control words {\tt~POINTS, GRIDTYPE, FILTER~} 
switches the input regime from scalar to the vector mode.
An order and case of the control words in a line are inessential.
Once defined, each control variable value is valid until it is redefined
by another control string.
If the six leading characters in an input string match any of the
control words then the following rules are applied: 
\begin{description}
\vspace*{-1ex}
\item[--] an order of control words in a control line is arbitrary.
\vspace*{-1ex}
\item[--] the word \twd{NAMEXP} must be present in the line, 
	and followed by a name from Table~5.2 or 5.4,
\vspace*{-1ex}
\item[--] the words \twd{NTIMES}, \twd{GRIDTYPE} and 
   \twd{POINTS} (if present) should be followed by an integer number,
\vspace*{-1ex}
\item[--] the words \twd{FACTOR} and \twd{FILTER} 
   (if present) should be followed by a real number,
\vspace*{-1ex}
\item[--] the words \twd{NTIMES, FACTOR} and \twd{FILTER} can be
omitted, then they acquire default values 1, 1. and 0.001, respectively.
\end{description}
\vspace*{-1ex}\noindent
A number after the word \twd{NTIMES} defines the number of time slices
in the input.  
These time values should always follow the control line.
A real number after the word \twd{FACTOR} defines the multiplication
(e.g. unit conversion) factor between the input data and in the Astra
variable.  
In the free format, no restrictions apart from those required by the
Fortran syntax are imposed. 

\paragraph{Free format for reading scalars.}
In the scalar variable input mode each control line should be followed
by a group of \twd{2*NTIMES} real numbers.
First \twd{NTIMES} of them give time and the rest give variable values.
For instance, the previous example \req{IPLX} can be written as
\smallskip\\
\dl{NAMEXP~IPL~~~~NTIMES~4~~~~FACTOR~1.E-6}
\dl{~.0~~~~~~~~1.~~~~~~~~2~~~~~~~~2.5}
\dl{.1e6~~~~~1.E+6~~~~~1.E6~~~~~~2.E5}
\vskip 10pt

The next example sets a box-like time dependence for the quantity
\twd{QECRX} 
\smallskip\\
\dl{NTIMES 4~~~~~NAMEXP QECR}
\dl{0.299~~~~~0.3~~~~~~1.5~~~~~~1.501}
\dl{0.0~~~~~~~1.0~~~~~~1.0~~~~~~0.0}
\vspace*{1ex}
It switches on from zero to 1 MW at $t=0.3$s and switches off at
$t=1.5$s. 

\paragraph{Free format for reading vector variables.}

The control words {\tt~POINTS, GRIDTYPE, FILTER~} switching the input
regime to the vector mode have the following meaning
\begin{description}
\vspace*{-1ex}
\item[--]\twd{POINTS} is a dimension of the input vector,
\vspace*{-1ex}
\item[--]\twd{GRIDTYPE} is an integer number specifying the radial
variable, 
\vspace*{-1ex}\item[--]
\twd{FILTER} is a real number determining a method of data transfer.
\end{description}
%\noindent
All control variables can be redefined after any input group as many
times as needed. 
Otherwise they extend their meaning to all subsequent groups with
one exception. 
The grid type and size, as determined by the variables \twd{GRIDTYPE}
and \twd{POINTS}, are not applied to the input from a U-file. 
The reason is that the input grid is defined independently in each
U-file and this definition is effective inside this U-file only.
When processing the U-file is over the variables \twd{GRIDTYPE} and 
\twd{POINTS} retrieve their values.
The action of the variable \twd{FILTER} is extended to all input formats.
The control words \twd{NAMEXP, NTIMES, FACTOR} and their 
values have the same significance and obey the same rules as for simple 
variables.
Input data should be split in strings no longer than 132 character each.

The variable \twd{POINTS} has no pre-defined value. 
It must appear in a control string at least once.
It sets the number of grid points to be used in the subsequent
input group. 
This grid has nothing to do with the grid used in the transport equation
solver. 
Therefore, a rule for data transfer from one grid to another is required.
Variables \twd{GRIDTYPE} and \twd{FILTER} define these rules.

The quantity \twd{GRIDTYPE} defines the radial variable and a type of 
the input grid.
In the current version of the code it can take the following values.
%\nopagebreak
\smallskip\\
\indent{\bf Table 5.5: Types of radial grid for input arrays}
\nopagebreak\\[10pt]\nopagebreak\indent
\begin{tabular}{|c|c|c|c|c|}
\hline
Value 	&  Radial coordinate & Type of the grid	& Units & Comment\\
\hline
0	& $a_j=\tw{AB*}\kappa_j$	& Equidistant & [m] & \\
1	& $a_j=\tw{ABC*}\kappa_j$	& Equidistant & [m] 
					& Default input mode \\
2	& $\rho_{N,j}=\kappa_j$	& Equidistant & [d/l] &	\\
3	& $\rho_{\psi,j}=\kappa_j$	& Equidistant & [d/l] & \\
4	& $\rho_{V,j}$		& Equidistant & [m] & Unimplemented \\
5	& $\psi_{N,j}$		& Equidistant & [d/l] & Unimplemented \\
6	& $\Phi_{j}$		& Equidistant & [Wb]  & Unimplemented \\
%6-9	& 			& Equidistant & [d/l] & Reserved  \\
10	& $a_j\in[0,\tw{AB}]$	& Arbitrary   & [m] &	 \\
11	& $a_j\in[0,\tw{ABC}]$	& Arbitrary   & [m] &	 \\
12	& $\rho_{N,j}\in[0,1]$	& Arbitrary   & [d/l] &	 \\
13	& $\rho_{\psi,j}\in[0,1]$& Arbitrary   & [d/l] & \\
14	& $\rho_{V,j}$		& Arbitrary   & [m] & 	 \\
15	& $\psi_{N,j}]$		& Arbitrary   & [d/l]& 	 \\
16	& $\Phi_{j}$		& Arbitrary   & [Wb]& 	 \\
18	& $\{r_j,z_0\}$		& Arbitrary 	& [m] & Vertical   chord \\
19	& $\{r_0,z_j\}$		& Arbitrary 	& [m] & Horizontal chord \\
20	& $\{r_j,z_j\}$		& Arbitrary 	& [m] & 2D grid \\
\hline
\end{tabular}
\\[15pt]
where $\kappa_j=\ds\frac{j-1}{\tw{POINTS}-1},\quad 1\leq j\leq\tw{POINTS}$ 
and
$$
\rho_N=\frac{\rho}{\rho_B}
,\qq
\rho_V=\sqrt\frac{V(\rho)-V(0)}{V(\rho_B)}\quad
,\qq
\psi_N=\frac{\psi(\rho)-\psi(0)}{\psi(\rho_B)-\psi(0)}
,\qq
\rho_\psi=\sqrt{\psi_N}
.
$$
When $\tw{GRIDTYPE}< 10,$ an equidistant grid is selected   
with respect to the radial coordinate used.
Note that equidistant grids on different radial variables
do not coincide with one another.

When $\tw{GRIDTYPE}\geq 10,$ the grid is arbitrary and is expected to be 
supplied by the user.
In this case, the grid can be extended also beyond the interval 
specified by the second column in Table 5.5
but all data beyond this interval will be ignored.
If the innermost point of the input grid does not coincide with 
the magnetic axis $\rho=0$ 
%or if the outermost point is inside the interval required 
then the input quantity at $\rho=0$ is set to the nearest value provided.
The same is done for the outer side.
In other words, every quantity is extended by a constant
outside the region where it is set by the input data 
if this region is smaller than the region of definition.

After that the input grid is mapped to the Astra transport grid by the
quadratic interpolation while the data are transferred to the transport
grid making use of the procedure described in Section~4.9.6 where
$\alpha=\tt FILTER.$ 

We consider now few examples. The first one reads: 
\\
\dl{******}
\dl{NAMEXP~~PRAD~~NTIMES~~2~~POINTS~~10~~GRIDTYPE~~1}
\dl{0.~~~0.03}
\dl{1.683E-02~~1.835E-02~~2.118E-02~~2.511E-02~~3.063E-02}
\dl{3.730E-02~~4.366E-02~~4.700E-02~~4.493E-02~~3.684E-02}
\dl{1.10~~0.95~~0.78~~0.63~~0.51~~0.45~~0.37~~0.25~~0.13~~0.04}
\dl{******}
\dl{POINTS~~11	GRIDTYPE~~10~~NAMEXP~~PRAD}
\dl{~.2}
\dl{~.0~~.05~~.1~~.15~~.2~~.25~~.3~~.35~~.4~~.45~~.5}
\dl{1.613903E-2~~1.683072E-2~~1.835108E-2~~2.118123E-2~~0.02511~~3.063486E-2}
\dl{3.730636E-2~~4.366064E-2~~4.700778E-2~~4.493699E-2~~3.684750E-2}
\dl{******}
Here stars separate successive input groups and are ignored by the code. 
The first control line declares that the input data should be written in
the array with the name \tw{PRADX} on the radial grid of 10 points which
is equidistant in the variable~$a$ on the segment $[0,\tw{ABC}].$ 
Two time slices ($t=0$s and $t=0.03$s) given in the next line 
are followed with 20 numbers which give the function values.
The second control line says that the input should be continued
using a new grid of 11 nodes. Then the time $t=0.2$s is given, 
then 11 data for $\{a_j\}$ and 11 values of the function.
Note that although the control pair \twd{"NTIMES~1"} is omitted the
corresponding time data must be present.

The second example:
\\
\dl{********************************* CAR3X *********************************}
\dl{>>>>	Data are given along a vertical chord }
\dl{POINTS~~10~~GRIDTYPE~~18~~NAMEXP~~CAR3}
\dl{~0.3}
\dl{~1.5}
\dl{-0.81 -0.64 -0.47 -0.30 -0.13~~0.04~~0.21~~0.38~~0.55~~0.72}
\dl{~0.04~~0.31~~0.56~~0.78~~0.93~~0.96~~0.87~~0.68~~0.44~~0.17}
\dl{********************************* CAR1X *********************************}
\dl{>>>>	Data are given along a downshifted horizontal chord }
\dl{POINTS~~11~~GRIDTYPE~~19~~NAMEXP~~CAR1}
\dl{~0.3}
\dl{-0.25}
\dl{~1.10~~1.20~~1.30~~1.40~~1.50~~1.60~~1.70~~1.80~~1.90~~2.00~~2.10}
\dl{-0.09~~0.26~~0.53~~0.72~~0.83~~0.87~~0.83~~0.72~~0.53~~0.26 -0.09}
\dl{********************************* CAR2X *********************************}
\dl{>>>>	Data are given along an inclined chord }
\dl{NAMEXP~~CAR2X~~GRIDTYPE~~20}
\dl{~0.3}
\dl{~1.10~~1.20~~1.30~~1.40~~1.50~~1.60~~1.70~~1.80~~1.90~~2.00~~2.10}
\dl{-0.30 -0.28 -0.27 -0.25 -0.24 -0.23 -0.21 -0.20 -0.18 -0.17 -0.15}
\dl{-0.12~~0.24~~0.51~~0.72~~0.84~~0.89~~0.87~~0.76~~0.58~~0.31 -0.03}
\dl{*************************************************************************}
defines three different quantities at the same time $t=$0.3 s 
which is specified by the first number after each control line.
In the first input group, \tw{GRIDTYPE=18}, therefore the time is
followed by the distance to the torus axis $r_0=1.5$m, then $\{z_j\}$
(also in meters) and $\{f_j\}$ are given, where $1\leq j\leq 10.$
The second control line switches the name, the grid type and size.  
It determines the horizontal chord with $z_0=-0.25$m, $\{r_j\}$ (in meters) 
and then gives the function values $\{f_j\}$
for $1\leq j\leq 11.$
The third group inherits \twd{POINTS=11} and after the time value gives
33 numbers for $\{r_j\},$ $\{z_j\}$ and $\{f_j\},$ with $1\leq j\leq 11.$

\subsubsection{Input format for the plasma boundary.}
When the three-moment equilibrium solver $(\tw{NEQUIL}\leq41)$ is
used the plasma boundary should be described by a set of scalar
variables \twd{RTOR, ABC, SHIFT, ELONG, TRIANG, UPDWN}.
This boundary setting can also be used for more complicated
equilibrium solvers, however, in the latter case, a more complicated
boundary shape is usually necessary.

A special input format is provided for setting the plasma boundary of
arbitrary shape.
The format is similar to the free format discussed above.
A new keyword \twd{BND} (or \tw{BNDX}) is introduced which declares that
the subsequent group of numbers will describe the plasma boundary.
One should beware that the boundary input switches the input mode from
scalar to vector type.

Plasma boundary input is illustrated with an example:
\\[1ex]
\dl{NAMEXP~~BND~~NTIMES~~2~~POINTS~~8}
\dl{0.263~~~0.313}
\dl{2.7750~~2.7000~~1.7631~~1.7695}
\dl{2.6283~~2.6371~-1.5053~-1.5104}
\dl{3.8605~~3.8605~~0.2516~~0.2516}
\dl{1.9041~~1.9167~~0.2516~~0.2516}
\dl{3.3000~~3.2250~~1.5363~~1.5932}
\dl{2.2748~~2.2500~~1.5172~~1.5014}
\dl{2.2515~~2.2604~-1.0141~-1.0141}
\dl{3.0728~~3.0750~-1.1406~-1.1526}
\\[1ex]
Here the time evolving plasma boundary is given by pairs of coordinates
$\{r_j,z_j\}$ where ${1\leq j\leq\tt POINTS=8.}$
First two time values $t_1=0.263$s and $t_2=0.313$s are given.
Then follow $2\times\tt NTIMES\times POINTS$ coordinates  
$r_1(t_1),\,r_1(t_2),\,z_1(t_1),\,z_1(t_2),\,r_2(t_1),\,r_2(t_2),\,
\dots\,,$\linebreak 
$r_8(t_1),\,r_8(t_2),\,z_8(t_1),\,z_8(t_2).$
It is understood that the boundary coordinates are linearly interpolated
if $t_1<t<t_2$ and do not evolve otherwise.

In the current version of the code all data for the boundary input
should be given in one group (multiple appearance of the combination 
\twd{"NAMEXP BND"} is not allowed) and the the total number of data in
this group $(2\times\tw{POINTS}+1)\times\tw{NTIMES}$ cannot exceed 25000.

\vskip 3ex
In conclusion, we outline a few {\bf common rules for a data file}:
\begin{description}
\vspace*{-1ex}
\item[--] The first two lines should be used for a device/shot 
identification and cannot contain a variable input.
\vspace*{-1ex}
\item[--] No meaningful (non-commentary) lines can use more than 132
positions. 
\vspace*{-1ex}
\item[--] If different time slices for the same quantity follow one
another then the time sequence should be put in increasing order.
\vspace*{-1ex}
\item[--] A sequence of records for every input quantity should be
contiguous, i.e. records for one quantity cannot alternate with records
for another. 
\vspace*{-1ex}
\item[--] At first, all simple variables should be determined. 
\vspace*{-1ex}
\item[--] A control line containing one of the words 
\twd{POINTS, GRIDTYPE, FILTER} switches the input to the vector mode.
After such a line radial profiles can be defined only. 
\vspace*{-1ex}
\item[--] A total number of input values for all simple variables and
all time slices cannot exceed a value of the parameter \twd{NTVAR} which
is set during the code installation. 
An actual value of the parameter is defined in the file
\twd{AWD/for/parameter.inc} (by default \twd{NTVAR=2000}).
For radial profiles the same limitation is given by the parameter
\twd{NTARR} (default value 5000).
\vspace*{-1ex}
\item[--] If the control word \tw{END} is encountered then the
subsequent contents of the data file is ignored. 
Otherwise, all records until the end of the file will be processed.
\end{description}

%\vfill\eject
%\section{Operation guide}
%\pagebreak
\subsection{Brief operation guide}

\subsubsection{Installing the code}

To install the Astra code the user is supplied with two files. 
One of them is an archived file containing all source files for the code.
Another is a script file \twd{"InstallAstra"}. 
Both files should be put in the same directory and then the code is
installed by running the script file. 
The script provides many options and can accept several control
parameters. 
For instance, there are options to install the code for multiple users
with a shared kernel, options which specify names for the
directory to put the code in and for the kernel to be linked to. 
Only a simplest option is considered here.

Running the script
\\ \indent\twd{InstallAstra}\\
without any control parameters should install the Astra kernel with the
name \twd{"Astra"} and the Astra user directory with the name
\twd{"astra"}. 
However, for an unknown platform some additional adjustments could be
requested. 
If an error is encountered a message is printed.
Then after a correction \twd{InstallAstra} should be run repeatedly.
Usually, after the Astra code is successfully installed it offers to set
two aliases 
\\ \indent
$\twd{Astra}=\twd{AWD/.exe/astra}$
\\
and
\\ \indent
$\twd{Review}=\twd{AWD/.exe/view}$
\\
where \twd{AWD} is replaced by the actual Astra working directory
name. 
Both alias names are arbitrary but for convenience these
shorthands will be used in what follows. 

\subsubsection{Starting the code}

After the Astra code is installed a syntax of the starting command and
the main contents of this section can be displayed by the command 
\\{\indent
\tt Astra help
}
\\
or
\\{\indent
\tt Astra -h
}
\\
The code is started by the command
\\{\indent
\tt Astra~[data\us file\us name]~[model\us name]~[start\us time]
[end\us time]\\
\hspace*{4.9em}
[post\us view\us file\us name]~[test]
}\\
which can take up to 6 parameters separated by spaces.
Any of them is optional and can be omitted. 
However, a meaning of every parameter is determined by its ordinal
position in the command line, therefore, if the dropped parameter is not
the last one in the line it should be replaced with a comma. 
Significance of the parameters should be obvious:

{\tt data\us file\us name~} is the data file name in the directory
\twd{AWD/exp/}, 

{\tt model\us name~~~~~} is the model name in the directory
\twd{AWD/equ/}, 

{\tt start\us time~~~~~} is the initial time of the simulation.\\
If one of these parameters is not defined then the value from the
previous run is used.
For instance, the command
\\{\indent
\tt Astra readme showdata
}\\
starts Astra run of a model described by the file
\twd{showdata} (full name \twd{AWD/equ/showdata}) with initial data
from \twd{readme} and initial time $t_0=0$ (default value).
To repeat the same run one needs to enter just
\\{\indent
\tt Astra
}\\
The command
\\{\indent
\tt Astra , analysis
}\\
can then be used to run the model \twd{analysis} with the same
data file \twd{readme}.

Remaining parameters:
 
{\tt end\us time~~~~~~~~~~~~} is the time when the simulation has
to be stopped, 

{\tt post\us view\us file\us name~} is the name for a file where the
results will be saved.\\
These two parameters are of use when interactive control is not required
and the code is run in the background. 
Then the 5-th parameter can be used for a post-run viewing the
results of modelling as described in Section~5.5.4. 
Normally, Astra writes all results of calculations in a file so that,
after the run, one can inspect the results of simulation as if they are
watched interactively.
These results are stored in a special post-view file and each new run
destroys the previous one.
Moreover, if several Astra processes run simultaneously they all will
write into the same file and corrupt it. 
In order to avoid this corruption during few simultaneous Astra runs 
the user has to employ the 5-th input parameter and submit the output of
different Astra processes to different files.

Finally, if the 6-th parameter \twd{"test"} or \twd{"TEST"} is set,
then the code execution is suspended after the all Fortran source files
are created (after Step~2 in the flow diagram on page \pageref{flow}). 
At this moment, the user can manually introduce any code modifications
which are not foreseen in the Astra compiler and resume execution of the
modified code. 

A superseding version of this command reads
\\{\indent
\tt Astra [-htVM] [-v data\us file\us name] [-m model\us name]
[-s start\us time] \\
\hspace*{4.9em}
[-e end\us time] [-r post\us view\us file\us name]
}\\
The command 
\\{\indent
\tt Astra -t
}\\
is equivalent to
\\{\indent
\tt Astra~,~,~,~,~,~test
}\\
The commands
\\{\indent
\tt Astra -V
}\\
and
\\{\indent
\tt Astra -M
}\\
list all available data files and models, respectively.

\subsubsection{Run time control}

When the Astra code is started it sets up an interactive mode.
In this mode, the user can control the simulation flow using menu boxes
displayed in the main window. 
If a menu box is selected with a mouse (any button) then either a
message is printed about the fulfilled action or the run is
suspended and a dialog window appears.
After carrying out required changes and exiting the dialog window the
run continues. 

The key \twd{<Esc>} always means exit from the active window. 
For a {\bf dialog window} pressing \twd{<Esc>} applies all changes,
closes the dialog window and returns cursor and focus to the main window.
When pressed in the {\bf main window \twd{<Esc>} stops} the code.
A dialog window displays a set of boxes containing numbers and words.
Only one box shown in bold (highlighted) is active at a time. 
The user can select a box either with the mouse or with the keyboard
arrows. 
Data in the selected box can be modified in order to change the run flow
as required. 
In a dialog window, useful is the key \twd{`\symbol{63}`} which prints
an information about the quantity in the highlighted box. 

All menu options of the main window are duplicated by keyboard keys
(letters are case insensitive). 
Most of menu items do not require much explanations therefore we
restrict this section to description of the main features of interactive
control. 
Instead of boring reading we encourage the user to test different options
and find the most suitable mode of operation. 

\paragraph{On-line help and information.}
The list of operational keys can be obtained by pressing \twd{`H`} or
\twd{`\symbol{63}`} or by selecting the menu box \twd{"Help"}. 
The key \twd{`L`} types the results of model analysis by the Astra
compiler. 
The current radial grid is explained by \twd{`X`}.
Correspondence between the menu options and the hot keys of this group
is shown in the table below.
\bigskip\\ 
\cl{\tt\indent
\begin{tabular}{| c | c | c | c |}\hline
\bf Menu option & Help		& Type model 	& What X-axis? 	\\ \hline
\bf Hot key 	& H {\rm or} ?	& L		& X		\\ \hline
\end{tabular}
}

\paragraph{Execution control} is performed by using the following menu
options: 
\bigskip\\ 
\cl{\tt\indent
\begin{tabular}{| c | c | c | c | c |}\hline
\bf Menu option & Variables & Constants & ~Grids~ & Save tuning \\ \hline
\bf Hot key 	& V	& C	& D	& I		\\ \hline
\end{tabular}
}
\\[2ex]
The keys \twd{`V`}, \twd{`D`} and \twd{`C`} invoke dialog windows which 
allow to change Astra variables listed in Tables~4.9--4.11, 4.13--4.16
and 4.17, respectively.
The key \twd{`I`} stores the current values of all these variables in a
file with the extension \twd{.log} as described in Section~5.3.2.

The set of variables accessible via a key \twd{`D`} includes 
also a set of control parameters defined by the command line \req{sbr}.
In this dialog window, one can change the control times for calling every
subroutine. 
There is an operation which can be done by the keyboard only.
It is usage of a control key as described in Section~5.24.
Pressing \twd{<Ctrl>} together with the predefined in the model key one
can call the corresponding subroutine out of order.
If such a combination matching one of the subroutine calling symbols is
pressed when code is in waiting mode the one time step is made.

Waiting and run modes are two alternative states of the code. 
Two keys \twd{<Space>} and \twd{<Return>} are reciprocal and set the
code in waiting and run modes, respectively.  
\bigskip\\
\cl{\tt
\begin{tabular}{| c | c | c | c || c |}\hline
\bf Menu option & Run   & Step    & Quit & Calling subroutine \\ \hline
\bf Hot key 	& <Return> & <Space> & <Esc> & <Ctrl><Letter> 	\\ \hline
\end{tabular}
}
\\[2ex]
The run mode is the default one when the calculations are performed and
the results presented on the screen are periodically renewed.
In the waiting mode, code does nothing although all keys are operable as
in the run mode.
In this mode, it is possible to advance calculation step by step
pressing \twd{<Space>} repeatedly.
Additionally, this mode can be used for mapping different radial
coordinates to one another because the cursor position is digitized and
displayed showing several flux labels simultaneously. 

\paragraph{Presentation control.}
Numerous keys are available for changing presentation of simulation
results. 
First of all, we will distinguish 9 different presentation modes.
Every of them can be set up pressing the corresponding digit-key.
The first three modes plot radially dependent functions:
\bigskip\\
\cl{\tt
\begin{tabular}{| c | c | c | c | c | c | c |}\hline
\bf Menu option &16*f(\#) & 8*f(\#) & 8*f(psi)	\\ \hline
\bf Hot key 	& 1	 & 2	   & 3		\\ \hline
\end{tabular}
}
\\[2ex]
In the first line, the number in front of \twd{*} shows a number of
curves plotted simultaneously at one screen.
Instead of the mark \twd{\#} in the symbolic notation above the Astra
code shows the actual radial coordinate which is defined by the
parameter \twd{XOUT} of Table~4.13 and can be one~of: 
\begin{itemize}
\vspace*{-1ex}
\item[--] $a$ introduced by \req{3mom}, 
		${\tt XOUT}=0$ or ${\tt XOUT}=1,$ 
\vspace*{-1ex}
\item[--] $\rho$ introduced by \req{dfPhi}, ${\tt XOUT}=2,$
\vspace*{-1ex}
\item[--] $\psi$ introduced by \req{dfpsi}, ${\tt XOUT}=3.$
\end{itemize}
\vspace*{-1ex}
The parameter \twd{XOUT} appears in the dialog window \twd{`D`} under
the name \twd{Xaxis}.
It can be changed at any time and stored by pressing \twd{`I`}.

The first (default) mode which simultaneously presents 16 radial
functions was briefly considered in Section~5.2.6 and depicted there
schematically as 
\bigskip\\ 
\cl{\tt
\begin{tabular}{| l | l | l | l |}\hline
\bf\large 1	&\bf\large 2	&\bf\large 3	&\bf\large 4	\\
1 \& 9		&2 \& 10	&3 \& 11	&4 \& 12	\\ \hline
\bf\large 5	&\bf\large 6	&\bf\large 7	&\bf\large 8 	\\
5 \& 13		&6 \& 14	&7 \& 15	&8 \& 16	\\ \hline
\end{tabular}
}
\\[2ex]
where the large numbers numerate boxes containing two radial profiles
each. 
The small digits are the sequence numbers of these radial profiles as
they appear in a model file.
If the key \twd{`N`} is pressed then the next 16 curves will be
displayed as shown below.
\bigskip\\ 
\cl{\tt
\begin{tabular}{| l | l | l | l |}				\hline
\bf\large 9	&\bf\large 10	&\bf\large 11	&\bf\large 12	\\
17 \& 25	&18 \& 26	&19 \& 27	&20 \& 28	\\ \hline
\bf\large 13	&\bf\large 14	&\bf\large 15	&\bf\large 16 	\\
21 \& 29	&22 \& 30	&23 \& 31	&24 \& 32	\\ \hline
\end{tabular}
}
\\[2ex]
If \twd{`N`} is pressed several times so that the list of profiles
requested for the radial output is exhausted then the first screen
appears again. 
The scan back can be carried out by pressing~\twd{`B`}.
\\[2ex]
\cl{\tt
\begin{tabular}{| c | c | c | c | c || c |}	\hline
\bf Menu option & Colors & Backward & Next & Refresh&Appearance\\ \hline
\bf Hot key 	& A & B & N & R & . 			\\ \hline
\end{tabular}
}
\\[2ex]
These keys operate similarly in all graphic modes. 
Colors can be changed in a dialog window opened by pressing \twd{`A`}.
The key \twd{`R`} redraws the current screen if it was corrupted for
any reason.
The hot key \twd{`.`} (dot) is especially useful for black and white
monitors or for plotting data on black and white printers. 
It changes appearance of the curves either making them dashed or marking
them with different symbols.

There are few other possibilities to change the appearance of curves.
\bigskip\\
\cl{\tt
\begin{tabular}{| c | c | c | c | c | c | c | c | c |}\hline
\bf Menu option	& Select & Scales & Windows & Y-shift	\\ \hline
\bf Hot key 	& M 	 & S 	  & W 	    & Y  	\\ \hline
\end{tabular}
}
\\[2ex]
The user can change scales for every curve \tw{`S`}, Y-offsets \tw{`Y`}
and also redirect any curve in another box \tw{`W`}. 
The key \tw{`M`} (option \twd{"Select"}) in modes 1, 2, 3, 6, 7 combines
actions of all three keys \tw{`S`}, \tw{`Y`} and \tw{`W`}.

%\vskip3ex
The second mode (hot key \twd{`2`}) shows the same profiles in larger
scale.
Repeated pressing of the key \twd{`2`} switches between two different
frames for the curves. 
In one case, the zero-ordinate line is in the bottom of a window, in
another, in the middle of it.\footnote
{A similar feature is also available for modes 3, 4, 5 and 6.}
This can be used for plotting curves with alternating sign.
Schematically the second mode can be shown as
\smallskip\\
\cl{\tt
\begin{tabular}{| l | l |}\hline
\bf\large 1	&\bf\large 2			\\
		&				\\
1 \& 5 \& 9 \& 13	&2 \& 6 \& 10 \& 14	\\
		&				\\ \hline
\end{tabular}
}
\\[2ex]
It is obtained if the lower row of boxes of the first mode is overlapped
with the upper row.
Consequently, the screen obtained by pressing \twd{`N`} can be
represented as 
\bigskip\\
\cl{\tt
\begin{tabular}{| l | l |}\hline
\bf\large 3	&\bf\large 4			\\
		&				\\
3 \& 7 \& 11 \& 15	&4 \& 8 \& 12 \& 16	\\
		&				\\ \hline
\end{tabular}
}
\\[2ex]
In other respects, the mode 2 is similar to the mode 1.

The mode 3 is obsolete because it always uses the coordinate $\psi$ for
the X-axis and, therefore, it is a particular case of the mode 2.
It will be removed in future versions.

The next three modes allow to plot time dependences
\bigskip\\
\cl{\tt
\begin{tabular}{| c | c | c | c | c | c | c |}		\hline
\bf Menu option & 2*f(a,t) & 2*f(R,t) & 8*f(t)		\\ \hline
\bf Hot key 	& 4	   & 5	      & 6		\\ \hline
\end{tabular}
}
\\[2ex]
The mode 6 shows eight time dependent curves placed as
\bigskip\\
\cl{\tt
\begin{tabular}{| l |}\hline
\bf\large 1~~~~~~~~~~~~~~~~~~~~~~~~~~~~~~~~~~~~~~~~~~~\\	\\
1 \& 3 \& 5 \& 7					\\ \hline
\bf\large 2~~~~~~~~~~~~~~~~~~~~~~~~~~~~~~~~~~~~~~~~~~~\\	\\
2 \& 4 \& 6 \& 8					\\ \hline
\end{tabular}
}
\\[2ex]
The time axis is defined by two parameters \twd{TINIT} and \twd{TSCALE}
included in the dialog window \twd{`D`}.
The first parameter defines position of the origin and the second one
the total length of the scale.
If in process of long simulation the current time $t$ becomes greater
than \twd{TINIT}+\twd{TSCALE} then the current evolution of parameters
is not seen any more.
In this situation, one can interactively change the parameters
\twd{TINIT} or \twd{TSCALE} and make the evolution visible again but
most probably it will require a new correction quite soon.
Sometimes, a better solution would be to set \twd{TSCALE} negative.
Then $\rm abs\,(\tw{TSCALE})$ will have the same sense as before but 
the $t$-axis will become floating.
In other words, \twd{TINIT} will be ignored and alignment of graphs will
be performed with respect to their right rather than left edge.

If the key \twd{`6`} is pressed in the sixth mode once more then the
frame will be changed to four boxes with two curves in each box
\bigskip\\
\cl{\tt
\begin{tabular}{| l |}\hline
\bf\large 1~~~~~~~~~~~~~~~~~~~~~~~~~~~~~~~~~~~~~~~~~~~\\
1 \& 5							\\ \hline
\bf\large 2~~~~~~~~~~~~~~~~~~~~~~~~~~~~~~~~~~~~~~~~~~~\\
2 \& 6							\\ \hline
\bf\large 3~~~~~~~~~~~~~~~~~~~~~~~~~~~~~~~~~~~~~~~~~~~\\
3 \& 7							\\ \hline
\bf\large 4~~~~~~~~~~~~~~~~~~~~~~~~~~~~~~~~~~~~~~~~~~~\\
4 \& 8							\\ \hline
\end{tabular}
}
\\[2ex]
The next pressing \twd{`6`} will put all eight curves in one box
%\bigskip\\ \cl{\tt\begin{tabular}{| l |}\hline
%\bf\large 1~~~~~~~~~~~~~~~~~~~~~~~~~~~~~~~~~~~~~~~~~~~\\ \\ \\	\\
%1 \& 2 \& 3 \& 4 \& 5 \& 6 \& 7 \& 8			\\ \hline
%\end{tabular}}\\[2ex]
and, finally, the forth pressing \twd{`6`} returns the drawing into
the original state. 

The modes 4 and 5 are similar and differ by the X-axis only.
Both show two radial profiles as functions of a coordinate in the
mid-plane. 
In the mode 4, this coordinate is $a$ (which is close to the minor
radius), in the mode 5, it is the major radius $R.$
In these two modes the radial profiles can be plotted for several time
slices simultaneously. 
When the mode is switched on all available radial profiles are shown
with shadow color.
Pressing the key \twd{`M`} one can select those profiles which should be
plotted and remove others.
In other respects, these two modes are similar to the mode 2.

The mode 7 can also be attributed to modes plotting time dependences.
It shows the same quantities as the mode 6 but as function one another.
In other words, it gives simulation trajectory in a phase space.
\bigskip\\
\cl{\tt
\begin{tabular}{| c | c | c | c | c |}\hline
\bf Menu option &Phase space&Equilibrium&User graph& No graphics\\ \hline
\bf Hot key     & 7	& 8	&9	&0	\\ \hline
\end{tabular}
}
\\[2ex] 
The mode 8 presents the plasma configuration (as a result of the
equilibrium solver) and some other informations as NBI geometry
or resonance position and so on. 
The mode `9` is reserved for plotting user's graphs.
It calls the subroutine \twd{AWD/sbr/mydraw.f} which can be edited by the
user and include any graphics not foreseen in the Astra code.
The key `0` suppresses the graphic output.

\paragraph{Output files.}
The user can save the figure which is currently on the screen as a
PostScript file either in portrait (\twd{`G`}) or in landscape
(\twd{`Q`}) orientation. 
The data can also be saved in three different text formats.
So the keys \twd{`F`} and \twd{`P`} results in writing ASCII files which
include either all radial profiles (if any of modes 1, 2, 3, 4 or 5 is
active) or all time dependences (in mode 6).
PostScript and ASCII files are written into directory \twd{AWD/dat/}.
Their names are generated automatically and reported in the command
window.
The same data as are stored in a file by pressing \twd{`F`} can be typed
to the screen with the key \twd{`T`}. 
\bigskip\\
\cl{\tt
\begin{tabular}{| c | c | c | c | c | c |}		\hline
\bf Menu option & Port\us PS & Land\us PS & Write data
				& Type data & U-files	\\ \hline
\bf Hot key 	& G & Q & F \& P & T & U		\\ \hline
\end{tabular}
}
\vskip2ex

It is also possible to save data in the U-file format.
When the key \twd{`U`} is pressed a dialog window is opened where a name
of the U-file to be stored should be given.
Note that a type of the U-file depends on the current mode of data
presentation. 
In the radial modes 1, 2 and 3, 1D U-files are created with the
radial variable which is currently on the screen.
Similarly, in the time mode 6, time dependent 1D U-files are written.
Finally, in the mode 4, a 2D U-file with both radial and time
dependences is created. 

\subsubsection{Post-run viewer}

Although the Astra code is basically designed for interactive work it
can also be run in background and the results may be inspected
afterwards. 
To enable this option the Astra code always saves the main results of
every run in a (post-view) file and provides a tool for its viewing.
A typical length of this file is few MB, therefore, if the file has not
been saved by an explicit command it will be overwritten by the next run
and lost. 
Therefore, the user should take care and issue the appropriate command if
results of the run should be saved. 
We remind that the simulation results can be submitted to a named file
directly by the Astra starting command (Section~5.5.1).

\paragraph{Review command syntax.}
The corresponding command in the full format reads
\\{\indent
\tt Review [-hl] [-d file\us name] [-s file\us name] [-r file\us name] 
		[file\us name]}
\\
The command 
\\{\tt\indent
Review -h
\hspace{40mm}{\rm or}\qq\qq
Review help}
\\
prints the short description of the Astra post-run viewer.
When called without parameters
\\{\tt\indent
Review}
\\
the command retrieves the most recent Astra run.
A name of the post-view file to be inspected can be given explicitly
\\{\tt\indent
Review file\us name}
\\
A list of all available post-view files can be obtained by
\\{\tt\indent
Review -l
\hspace{40mm}{\rm or}\qq\qq
Review @}
\\
The saving command is 
\\{\tt\indent
Review -s file\us name
\qq\qq{\rm ~or}\qq\qq
Review file\us name save}
\\
An existing post-view file can be removed with the command
\\{\tt\indent
Review -d file\us name
\qq\qq{\rm ~or}\qq\qq
Review file\us name del[ete]}

\paragraph{Operation control.}
The \twd{Review} command opens a graphic window which is very similar in
image and operation to the Astra run window. 
Similar to the run mode, there is a selection of menu boxes duplicated
with control keys which mostly coincide in both modes.
Moreover, some additional keys (arrows and \twd{<Home>,}
\twd{<End>,} \twd{<PgUp>,} \twd{<PgDn>}) are enabled. 
These keys serve to change the current viewing time (shown in the right
upper corner) thus selecting appropriate time slices for radial profile
plotting. 
The arrows advance the current time either in small ($\gets$ and $\to$)
or in large ($\uparrow$ and $\downarrow$) steps scanning the simulation
time up and down, respectively.
\twd{<Home>} (\twd{<End>}) moves the viewing time to the beginning (end)
of the record. 
Two additional keys \twd{<PgUp>} and \twd{<PgDn>} can be used if the
review file is too long and can be loaded by parts only.

On the other hand, some limitations for changing control parameters are
added because no calculation is made in the view mode. 
For instance, the keys \twd{`V`} and \twd{`C`} can be used for
inspecting the current values of control or plasma parameters which
cannot be changed in the view mode.
The full list of active keys is printed when the key \twd{`H`} is
pressed. 

\paragraph{Model retrieval.}

Transport modelling is controlled by a huge amount of plasma and tuning
parameters.
Tracking all changes in these parameters which are continuously
varying is not easy.
Therefore, a problem of reproducing a simulation run performed a while
ago could be quite messy.
To this end, Astra review file includes also a complete information
about the model and tuning parameters for every saved run.

The model file and the corresponding start $(\tt.log)$ file can be
recovered with the command
\\{\tt\indent
Review -r file\us name}
\\
Here \twd{file\us name} is the post-view file name.
The retrieved model and \twd{.log} files will be created in the Astra
working directory \twd{AWD/} under the names \twd{model.tmp} and
\twd{model.log}, respectively. 
Then the user will be prompted to save these files under unique names in
order not to destroy them by the subsequent operations.

Note, however, that the initial data file is not stored.
Therefore, if it was changed in the meantime then the exact reproduction 
could be not possible. 
A similar problem arises if the tuning parameters have been modified
during the run.
Although the history of parameter variation can be restored by a careful
examination of the post-view file there is no automatic procedure for
doing this. 
Moreover, this history is stored in a post-view file with some
discreteness, therefore, this information still is incomplete.

%\vfill\eject



